\section{清初其余账本事件锚点}
本节补齐清初账本中尚未导航的条目，使战争、行政、海疆、边地与社会后果的不同材料层级均可直接跳转。
\subsection{1650—1659}
\eventanchor{EQ-1653-ZHENG-XIAMEN-AFTERMATH}{「郑成功以厦门、金门为基地整…}顺治十年（1653年），「郑成功以厦门、金门为基地整合海上军政与贸易网络，持续攻击福建沿海清军据点」引发的善后、执行或地方回应构成可由不同材料追踪的后续节点。核心取证：《清世祖实录》顺治十年福建海防军报；《清史稿·郑成功传》；《厦门志》与郑氏文书辑录；相关地区的档案、志书、契约、报刊或对方材料。该项锚定的是后续追踪问题，影响范围须按地区与群体分别判断。来源保留：战后、改革后或条约后的记忆、地方记录和官方总结常具有事后选择性；后续影响须按地点、群体和时间分层，不作整体化推断。
\eventanchor{EQ-1654-ZHENG-ZHANGZHOU-PRELUDE}{「郑成功军围攻漳州等闽南城镇…}顺治十一年（1654年），「郑成功军围攻漳州等闽南城镇，清军守城与援军维持陆海据点」之前的条件、信息和争议已通过中央与地方材料显现，构成该事件可追踪的前史节点。核心取证：《清世祖实录》顺治十一年福建军报；《清史稿·郑成功传》；《漳州府志》与荷兰东印度公司台湾商馆日记；同时期地方档案、地方志、日记或对方材料应按具体地区补检。该项锚定的是前史条件与信息背景，不把条件直接写成唯一因果。来源保留：官修编年和地方材料的记录密度与立场并不相同；本行不将条件、传闻、呈报或后见解释直接等同为确定因果。
\eventanchor{EQ-1654-ZHENG-ZHANGZHOU-DECISION}{围绕「郑成功军围攻漳州等闽南…}顺治十一年（1654年），围绕「郑成功军围攻漳州等闽南城镇，清军守城与援军维持陆海据点」的命令、调度、照会或呈报形成独立的中央—地方行动文书链。核心取证：《清世祖实录》顺治十一年福建军报；《清史稿·郑成功传》；《漳州府志》与荷兰东印度公司台湾商馆日记。该项锚定的是命令或决策环节，命令抵达和结果仍须另外核验。来源保留：奏折、上谕、部议、军机档或条约可证明行政动作和机构立场，不自动证明所有地区、群体或对手按其意图行动。
\eventanchor{EQ-1654-ZHENG-ZHANGZHOU-AFTERMATH}{「郑成功军围攻漳州等闽南城镇…}顺治十一年（1654年），「郑成功军围攻漳州等闽南城镇，清军守城与援军维持陆海据点」引发的善后、执行或地方回应构成可由不同材料追踪的后续节点。核心取证：《清世祖实录》顺治十一年福建军报；《清史稿·郑成功传》；《漳州府志》与荷兰东印度公司台湾商馆日记；相关地区的档案、志书、契约、报刊或对方材料。该项锚定的是后续追踪问题，影响范围须按地区与群体分别判断。来源保留：战后、改革后或条约后的记忆、地方记录和官方总结常具有事后选择性；后续影响须按地点、群体和时间分层，不作整体化推断。
\eventanchor{EQ-1656-ZHENG-QUANZHOU-PRELUDE}{「郑成功军在泉州、同安一线进…}顺治十三年（1656年），「郑成功军在泉州、同安一线进攻并争夺闽南水陆通道」之前的条件、信息和争议已通过中央与地方材料显现，构成该事件可追踪的前史节点。核心取证：《清世祖实录》顺治十三年福建军报；《清史稿·郑成功传》；《泉州府志》及郑氏文书辑录；同时期地方档案、地方志、日记或对方材料应按具体地区补检。该项锚定的是前史条件与信息背景，不把条件直接写成唯一因果。来源保留：官修编年和地方材料的记录密度与立场并不相同；本行不将条件、传闻、呈报或后见解释直接等同为确定因果。
\eventanchor{EQ-1656-ZHENG-QUANZHOU-DECISION}{围绕「郑成功军在泉州、同安一…}顺治十三年（1656年），围绕「郑成功军在泉州、同安一线进攻并争夺闽南水陆通道」的命令、调度、照会或呈报形成独立的中央—地方行动文书链。核心取证：《清世祖实录》顺治十三年福建军报；《清史稿·郑成功传》；《泉州府志》及郑氏文书辑录。该项锚定的是命令或决策环节，命令抵达和结果仍须另外核验。来源保留：奏折、上谕、部议、军机档或条约可证明行政动作和机构立场，不自动证明所有地区、群体或对手按其意图行动。
\eventanchor{EQ-1656-ZHENG-QUANZHOU-AFTERMATH}{「郑成功军在泉州、同安一线进…}顺治十三年（1656年），「郑成功军在泉州、同安一线进攻并争夺闽南水陆通道」引发的善后、执行或地方回应构成可由不同材料追踪的后续节点。核心取证：《清世祖实录》顺治十三年福建军报；《清史稿·郑成功传》；《泉州府志》及郑氏文书辑录；相关地区的档案、志书、契约、报刊或对方材料。该项锚定的是后续追踪问题，影响范围须按地区与群体分别判断。来源保留：战后、改革后或条约后的记忆、地方记录和官方总结常具有事后选择性；后续影响须按地点、群体和时间分层，不作整体化推断。
\eventanchor{EQ-1657-YONGLI-GUIZHOU-PRELUDE}{「永历朝在贵州、云南的孙可望…}永历十一年／顺治十四年（1657年），「永历朝在贵州、云南的孙可望与李定国集团矛盾公开化，孙可望转降清廷」之前的条件、信息和争议已通过中央与地方材料显现，构成该事件可追踪的前史节点。核心取证：《清世祖实录》顺治十四年贵州军报；《清史稿·孙可望传》；《永历实录》与贵州地方志；同时期地方档案、地方志、日记或对方材料应按具体地区补检。该项锚定的是前史条件与信息背景，不把条件直接写成唯一因果。来源保留：官修编年和地方材料的记录密度与立场并不相同；本行不将条件、传闻、呈报或后见解释直接等同为确定因果。
\eventanchor{EQ-1657-YONGLI-GUIZHOU-DECISION}{围绕「永历朝在贵州、云南的孙…}永历十一年／顺治十四年（1657年），围绕「永历朝在贵州、云南的孙可望与李定国集团矛盾公开化，孙可望转降清廷」的命令、调度、照会或呈报形成独立的中央—地方行动文书链。核心取证：《清世祖实录》顺治十四年贵州军报；《清史稿·孙可望传》；《永历实录》与贵州地方志。该项锚定的是命令或决策环节，命令抵达和结果仍须另外核验。来源保留：奏折、上谕、部议、军机档或条约可证明行政动作和机构立场，不自动证明所有地区、群体或对手按其意图行动。
\eventanchor{EQ-1657-YONGLI-GUIZHOU-AFTERMATH}{「永历朝在贵州、云南的孙可望…}永历十一年／顺治十四年（1657年），「永历朝在贵州、云南的孙可望与李定国集团矛盾公开化，孙可望转降清廷」引发的善后、执行或地方回应构成可由不同材料追踪的后续节点。核心取证：《清世祖实录》顺治十四年贵州军报；《清史稿·孙可望传》；《永历实录》与贵州地方志；相关地区的档案、志书、契约、报刊或对方材料。该项锚定的是后续追踪问题，影响范围须按地区与群体分别判断。来源保留：战后、改革后或条约后的记忆、地方记录和官方总结常具有事后选择性；后续影响须按地点、群体和时间分层，不作整体化推断。
\eventanchor{EQ-1658-QING-YUNNAN-PRELUDE}{「清军由贵州、四川等方向进逼…}顺治十五年（1658年），「清军由贵州、四川等方向进逼云南，永历朝廷撤离昆明并向西南退却」之前的条件、信息和争议已通过中央与地方材料显现，构成该事件可追踪的前史节点。核心取证：《清世祖实录》顺治十五年云南军报；《清史稿·洪承畴传》；《云南通志》与永历朝奏疏辑录；同时期地方档案、地方志、日记或对方材料应按具体地区补检。该项锚定的是前史条件与信息背景，不把条件直接写成唯一因果。来源保留：官修编年和地方材料的记录密度与立场并不相同；本行不将条件、传闻、呈报或后见解释直接等同为确定因果。
\eventanchor{EQ-1658-QING-YUNNAN-DECISION}{围绕「清军由贵州、四川等方向…}顺治十五年（1658年），围绕「清军由贵州、四川等方向进逼云南，永历朝廷撤离昆明并向西南退却」的命令、调度、照会或呈报形成独立的中央—地方行动文书链。核心取证：《清世祖实录》顺治十五年云南军报；《清史稿·洪承畴传》；《云南通志》与永历朝奏疏辑录。该项锚定的是命令或决策环节，命令抵达和结果仍须另外核验。来源保留：奏折、上谕、部议、军机档或条约可证明行政动作和机构立场，不自动证明所有地区、群体或对手按其意图行动。
\eventanchor{EQ-1658-QING-YUNNAN-AFTERMATH}{「清军由贵州、四川等方向进逼…}顺治十五年（1658年），「清军由贵州、四川等方向进逼云南，永历朝廷撤离昆明并向西南退却」引发的善后、执行或地方回应构成可由不同材料追踪的后续节点。核心取证：《清世祖实录》顺治十五年云南军报；《清史稿·洪承畴传》；《云南通志》与永历朝奏疏辑录；相关地区的档案、志书、契约、报刊或对方材料。该项锚定的是后续追踪问题，影响范围须按地区与群体分别判断。来源保留：战后、改革后或条约后的记忆、地方记录和官方总结常具有事后选择性；后续影响须按地点、群体和时间分层，不作整体化推断。
\eventanchor{EQ-1659-NANJING-EXPEDITION-PRELUDE}{「郑成功北上围攻南京，未能突…}顺治十六年七月至八月（1659年），「郑成功北上围攻南京，未能突破城防后撤回海上基地」之前的条件、信息和争议已通过中央与地方材料显现，构成该事件可追踪的前史节点。核心取证：《清世祖实录》顺治十六年七月至八月；《清史稿·郑成功传》；《延平王户官杨英从征实录》与荷兰东印度公司海上报告；同时期地方档案、地方志、日记或对方材料应按具体地区补检。该项锚定的是前史条件与信息背景，不把条件直接写成唯一因果。来源保留：官修编年和地方材料的记录密度与立场并不相同；本行不将条件、传闻、呈报或后见解释直接等同为确定因果。
\eventanchor{EQ-1659-NANJING-EXPEDITION-DECISION}{围绕「郑成功北上围攻南京，未…}顺治十六年七月至八月（1659年），围绕「郑成功北上围攻南京，未能突破城防后撤回海上基地」的命令、调度、照会或呈报形成独立的中央—地方行动文书链。核心取证：《清世祖实录》顺治十六年七月至八月；《清史稿·郑成功传》；《延平王户官杨英从征实录》与荷兰东印度公司海上报告。该项锚定的是命令或决策环节，命令抵达和结果仍须另外核验。来源保留：奏折、上谕、部议、军机档或条约可证明行政动作和机构立场，不自动证明所有地区、群体或对手按其意图行动。
\eventanchor{EQ-1659-NANJING-EXPEDITION-AFTERMATH}{「郑成功北上围攻南京，未能突…}顺治十六年七月至八月（1659年），「郑成功北上围攻南京，未能突破城防后撤回海上基地」引发的善后、执行或地方回应构成可由不同材料追踪的后续节点。核心取证：《清世祖实录》顺治十六年七月至八月；《清史稿·郑成功传》；《延平王户官杨英从征实录》与荷兰东印度公司海上报告；相关地区的档案、志书、契约、报刊或对方材料。该项锚定的是后续追踪问题，影响范围须按地区与群体分别判断。来源保留：战后、改革后或条约后的记忆、地方记录和官方总结常具有事后选择性；后续影响须按地点、群体和时间分层，不作整体化推断。
\eventanchor{EQ-1659-YONGLI-MYANMAR-PRELUDE}{「永历帝一行进入缅甸，西南战…}永历十三年／顺治十六年末（1659年），「永历帝一行进入缅甸，西南战争扩展为跨境避难和外交问题」之前的条件、信息和争议已通过中央与地方材料显现，构成该事件可追踪的前史节点。核心取证：《清世祖实录》顺治十六年云南军报；《清史稿·永历帝本纪》；《琉璃宫史》缅文传统及欧洲旅行记译注；同时期地方档案、地方志、日记或对方材料应按具体地区补检。该项锚定的是前史条件与信息背景，不把条件直接写成唯一因果。来源保留：官修编年和地方材料的记录密度与立场并不相同；本行不将条件、传闻、呈报或后见解释直接等同为确定因果。
\eventanchor{EQ-1659-YONGLI-MYANMAR-DECISION}{围绕「永历帝一行进入缅甸，西…}永历十三年／顺治十六年末（1659年），围绕「永历帝一行进入缅甸，西南战争扩展为跨境避难和外交问题」的命令、调度、照会或呈报形成独立的中央—地方行动文书链。核心取证：《清世祖实录》顺治十六年云南军报；《清史稿·永历帝本纪》；《琉璃宫史》缅文传统及欧洲旅行记译注。该项锚定的是命令或决策环节，命令抵达和结果仍须另外核验。来源保留：奏折、上谕、部议、军机档或条约可证明行政动作和机构立场，不自动证明所有地区、群体或对手按其意图行动。
\eventanchor{EQ-1659-YONGLI-MYANMAR-AFTERMATH}{「永历帝一行进入缅甸，西南战…}永历十三年／顺治十六年末（1659年），「永历帝一行进入缅甸，西南战争扩展为跨境避难和外交问题」引发的善后、执行或地方回应构成可由不同材料追踪的后续节点。核心取证：《清世祖实录》顺治十六年云南军报；《清史稿·永历帝本纪》；《琉璃宫史》缅文传统及欧洲旅行记译注；相关地区的档案、志书、契约、报刊或对方材料。该项锚定的是后续追踪问题，影响范围须按地区与群体分别判断。来源保留：战后、改革后或条约后的记忆、地方记录和官方总结常具有事后选择性；后续影响须按地点、群体和时间分层，不作整体化推断。
\subsection{1660—1669}
\eventanchor{EQ-1660-QING-YUNNAN-CONSOLIDATION-PRELUDE}{「清军继续经营云南城镇与驿路…}顺治十七年（1660年），「清军继续经营云南城镇与驿路，永历朝残部向中缅边地退缩」之前的条件、信息和争议已通过中央与地方材料显现，构成该事件可追踪的前史节点。核心取证：《清世祖实录》顺治十七年云南军报；《清史稿·吴三桂传》；内阁大库顺治十七年云南粮饷题本与《云南通志》；同时期地方档案、地方志、日记或对方材料应按具体地区补检。该项锚定的是前史条件与信息背景，不把条件直接写成唯一因果。来源保留：官修编年和地方材料的记录密度与立场并不相同；本行不将条件、传闻、呈报或后见解释直接等同为确定因果。
\eventanchor{EQ-1660-QING-YUNNAN-CONSOLIDATION-DECISION}{围绕「清军继续经营云南城镇与…}顺治十七年（1660年），围绕「清军继续经营云南城镇与驿路，永历朝残部向中缅边地退缩」的命令、调度、照会或呈报形成独立的中央—地方行动文书链。核心取证：《清世祖实录》顺治十七年云南军报；《清史稿·吴三桂传》；内阁大库顺治十七年云南粮饷题本与《云南通志》。该项锚定的是命令或决策环节，命令抵达和结果仍须另外核验。来源保留：奏折、上谕、部议、军机档或条约可证明行政动作和机构立场，不自动证明所有地区、群体或对手按其意图行动。
\eventanchor{EQ-1660-QING-YUNNAN-CONSOLIDATION-AFTERMATH}{「清军继续经营云南城镇与驿路…}顺治十七年（1660年），「清军继续经营云南城镇与驿路，永历朝残部向中缅边地退缩」引发的善后、执行或地方回应构成可由不同材料追踪的后续节点。核心取证：《清世祖实录》顺治十七年云南军报；《清史稿·吴三桂传》；内阁大库顺治十七年云南粮饷题本与《云南通志》；相关地区的档案、志书、契约、报刊或对方材料。该项锚定的是后续追踪问题，影响范围须按地区与群体分别判断。来源保留：战后、改革后或条约后的记忆、地方记录和官方总结常具有事后选择性；后续影响须按地点、群体和时间分层，不作整体化推断。
\eventanchor{EQ-1661-KANGXI-ACCESSION-PRELUDE}{「顺治帝去世，八岁玄烨即位为…}顺治十八年正月（1661年），「顺治帝去世，八岁玄烨即位为康熙帝，索尼、苏克萨哈、遏必隆、鳌拜辅政」之前的条件、信息和争议已通过中央与地方材料显现，构成该事件可追踪的前史节点。核心取证：《清世祖实录》顺治十八年正月；《清史稿·圣祖本纪》；《康熙起居注》与内国史院遗诏档；同时期地方档案、地方志、日记或对方材料应按具体地区补检。该项锚定的是前史条件与信息背景，不把条件直接写成唯一因果。来源保留：官修编年和地方材料的记录密度与立场并不相同；本行不将条件、传闻、呈报或后见解释直接等同为确定因果。
\eventanchor{EQ-1661-KANGXI-ACCESSION-DECISION}{围绕「顺治帝去世，八岁玄烨即…}顺治十八年正月（1661年），围绕「顺治帝去世，八岁玄烨即位为康熙帝，索尼、苏克萨哈、遏必隆、鳌拜辅政」的命令、调度、照会或呈报形成独立的中央—地方行动文书链。核心取证：《清世祖实录》顺治十八年正月；《清史稿·圣祖本纪》；《康熙起居注》与内国史院遗诏档。该项锚定的是命令或决策环节，命令抵达和结果仍须另外核验。来源保留：奏折、上谕、部议、军机档或条约可证明行政动作和机构立场，不自动证明所有地区、群体或对手按其意图行动。
\eventanchor{EQ-1661-KANGXI-ACCESSION-AFTERMATH}{「顺治帝去世，八岁玄烨即位为…}顺治十八年正月（1661年），「顺治帝去世，八岁玄烨即位为康熙帝，索尼、苏克萨哈、遏必隆、鳌拜辅政」引发的善后、执行或地方回应构成可由不同材料追踪的后续节点。核心取证：《清世祖实录》顺治十八年正月；《清史稿·圣祖本纪》；《康熙起居注》与内国史院遗诏档；相关地区的档案、志书、契约、报刊或对方材料。该项锚定的是后续追踪问题，影响范围须按地区与群体分别判断。来源保留：战后、改革后或条约后的记忆、地方记录和官方总结常具有事后选择性；后续影响须按地点、群体和时间分层，不作整体化推断。
\eventanchor{EQ-1661-COAST-EVACUATION-PRELUDE}{「清廷下令福建广东沿海迁界，…}顺治十八年（1661年），「清廷下令福建广东沿海迁界，试图切断郑氏的粮食、人口和贸易联系」之前的条件、信息和争议已通过中央与地方材料显现，构成该事件可追踪的前史节点。核心取证：《清世祖实录》顺治十八年迁界谕旨；《清史稿·食货志》；《福建通志》《广东通志》及沿海县志迁界条；同时期地方档案、地方志、日记或对方材料应按具体地区补检。该项锚定的是前史条件与信息背景，不把条件直接写成唯一因果。来源保留：官修编年和地方材料的记录密度与立场并不相同；本行不将条件、传闻、呈报或后见解释直接等同为确定因果。
\eventanchor{EQ-1661-COAST-EVACUATION-DECISION}{围绕「清廷下令福建广东沿海迁…}顺治十八年（1661年），围绕「清廷下令福建广东沿海迁界，试图切断郑氏的粮食、人口和贸易联系」的命令、调度、照会或呈报形成独立的中央—地方行动文书链。核心取证：《清世祖实录》顺治十八年迁界谕旨；《清史稿·食货志》；《福建通志》《广东通志》及沿海县志迁界条。该项锚定的是命令或决策环节，命令抵达和结果仍须另外核验。来源保留：奏折、上谕、部议、军机档或条约可证明行政动作和机构立场，不自动证明所有地区、群体或对手按其意图行动。
\eventanchor{EQ-1661-COAST-EVACUATION-AFTERMATH}{「清廷下令福建广东沿海迁界，…}顺治十八年（1661年），「清廷下令福建广东沿海迁界，试图切断郑氏的粮食、人口和贸易联系」引发的善后、执行或地方回应构成可由不同材料追踪的后续节点。核心取证：《清世祖实录》顺治十八年迁界谕旨；《清史稿·食货志》；《福建通志》《广东通志》及沿海县志迁界条；相关地区的档案、志书、契约、报刊或对方材料。该项锚定的是后续追踪问题，影响范围须按地区与群体分别判断。来源保留：战后、改革后或条约后的记忆、地方记录和官方总结常具有事后选择性；后续影响须按地点、群体和时间分层，不作整体化推断。
\eventanchor{EQ-1661-ZHENG-TAIWAN-LANDING-PRELUDE}{「郑成功军登陆台湾并围攻荷兰…}永历十五年／顺治十八年四月（1661年），「郑成功军登陆台湾并围攻荷兰东印度公司在热兰遮城的据点」之前的条件、信息和争议已通过中央与地方材料显现，构成该事件可追踪的前史节点。核心取证：《清世祖实录》顺治十八年海防军报；《清史稿·郑成功传》；《热兰遮城日志》荷文1661年四月至十二月条；同时期地方档案、地方志、日记或对方材料应按具体地区补检。该项锚定的是前史条件与信息背景，不把条件直接写成唯一因果。来源保留：官修编年和地方材料的记录密度与立场并不相同；本行不将条件、传闻、呈报或后见解释直接等同为确定因果。
\eventanchor{EQ-1661-ZHENG-TAIWAN-LANDING-DECISION}{围绕「郑成功军登陆台湾并围攻…}永历十五年／顺治十八年四月（1661年），围绕「郑成功军登陆台湾并围攻荷兰东印度公司在热兰遮城的据点」的命令、调度、照会或呈报形成独立的中央—地方行动文书链。核心取证：《清世祖实录》顺治十八年海防军报；《清史稿·郑成功传》；《热兰遮城日志》荷文1661年四月至十二月条。该项锚定的是命令或决策环节，命令抵达和结果仍须另外核验。来源保留：奏折、上谕、部议、军机档或条约可证明行政动作和机构立场，不自动证明所有地区、群体或对手按其意图行动。
\eventanchor{EQ-1661-ZHENG-TAIWAN-LANDING-AFTERMATH}{「郑成功军登陆台湾并围攻荷兰…}永历十五年／顺治十八年四月（1661年），「郑成功军登陆台湾并围攻荷兰东印度公司在热兰遮城的据点」引发的善后、执行或地方回应构成可由不同材料追踪的后续节点。核心取证：《清世祖实录》顺治十八年海防军报；《清史稿·郑成功传》；《热兰遮城日志》荷文1661年四月至十二月条；相关地区的档案、志书、契约、报刊或对方材料。该项锚定的是后续追踪问题，影响范围须按地区与群体分别判断。来源保留：战后、改革后或条约后的记忆、地方记录和官方总结常具有事后选择性；后续影响须按地点、群体和时间分层，不作整体化推断。
\eventanchor{EQ-1662-ZEELANDIA-PRELUDE}{「热兰遮城守军投降，荷兰东印…}永历十六年／康熙元年正月（1662年），「热兰遮城守军投降，荷兰东印度公司在台湾西南的统治终结」之前的条件、信息和争议已通过中央与地方材料显现，构成该事件可追踪的前史节点。核心取证：《清圣祖实录》康熙元年台湾相关军报；《清史稿·郑成功传》；《热兰遮城日志》荷文1662年正月及投降条约；同时期地方档案、地方志、日记或对方材料应按具体地区补检。该项锚定的是前史条件与信息背景，不把条件直接写成唯一因果。来源保留：官修编年和地方材料的记录密度与立场并不相同；本行不将条件、传闻、呈报或后见解释直接等同为确定因果。
\eventanchor{EQ-1662-ZEELANDIA-DECISION}{围绕「热兰遮城守军投降，荷兰…}永历十六年／康熙元年正月（1662年），围绕「热兰遮城守军投降，荷兰东印度公司在台湾西南的统治终结」的命令、调度、照会或呈报形成独立的中央—地方行动文书链。核心取证：《清圣祖实录》康熙元年台湾相关军报；《清史稿·郑成功传》；《热兰遮城日志》荷文1662年正月及投降条约。该项锚定的是命令或决策环节，命令抵达和结果仍须另外核验。来源保留：奏折、上谕、部议、军机档或条约可证明行政动作和机构立场，不自动证明所有地区、群体或对手按其意图行动。
\eventanchor{EQ-1662-ZEELANDIA-AFTERMATH}{「热兰遮城守军投降，荷兰东印…}永历十六年／康熙元年正月（1662年），「热兰遮城守军投降，荷兰东印度公司在台湾西南的统治终结」引发的善后、执行或地方回应构成可由不同材料追踪的后续节点。核心取证：《清圣祖实录》康熙元年台湾相关军报；《清史稿·郑成功传》；《热兰遮城日志》荷文1662年正月及投降条约；相关地区的档案、志书、契约、报刊或对方材料。该项锚定的是后续追踪问题，影响范围须按地区与群体分别判断。来源保留：战后、改革后或条约后的记忆、地方记录和官方总结常具有事后选择性；后续影响须按地点、群体和时间分层，不作整体化推断。
\eventanchor{EQ-1662-YONGLI-EXECUTION-PRELUDE}{「吴三桂军将自缅甸押回的永历…}康熙元年四月（1662年），「吴三桂军将自缅甸押回的永历帝朱由榔处死，南明朝廷名义中心结束」之前的条件、信息和争议已通过中央与地方材料显现，构成该事件可追踪的前史节点。核心取证：《清圣祖实录》康熙元年四月；《清史稿·永历帝本纪》；《琉璃宫史》缅文传统与《云南通志》；同时期地方档案、地方志、日记或对方材料应按具体地区补检。该项锚定的是前史条件与信息背景，不把条件直接写成唯一因果。来源保留：官修编年和地方材料的记录密度与立场并不相同；本行不将条件、传闻、呈报或后见解释直接等同为确定因果。
\eventanchor{EQ-1662-YONGLI-EXECUTION-DECISION}{围绕「吴三桂军将自缅甸押回的…}康熙元年四月（1662年），围绕「吴三桂军将自缅甸押回的永历帝朱由榔处死，南明朝廷名义中心结束」的命令、调度、照会或呈报形成独立的中央—地方行动文书链。核心取证：《清圣祖实录》康熙元年四月；《清史稿·永历帝本纪》；《琉璃宫史》缅文传统与《云南通志》。该项锚定的是命令或决策环节，命令抵达和结果仍须另外核验。来源保留：奏折、上谕、部议、军机档或条约可证明行政动作和机构立场，不自动证明所有地区、群体或对手按其意图行动。
\eventanchor{EQ-1662-YONGLI-EXECUTION-AFTERMATH}{「吴三桂军将自缅甸押回的永历…}康熙元年四月（1662年），「吴三桂军将自缅甸押回的永历帝朱由榔处死，南明朝廷名义中心结束」引发的善后、执行或地方回应构成可由不同材料追踪的后续节点。核心取证：《清圣祖实录》康熙元年四月；《清史稿·永历帝本纪》；《琉璃宫史》缅文传统与《云南通志》；相关地区的档案、志书、契约、报刊或对方材料。该项锚定的是后续追踪问题，影响范围须按地区与群体分别判断。来源保留：战后、改革后或条约后的记忆、地方记录和官方总结常具有事后选择性；后续影响须按地点、群体和时间分层，不作整体化推断。
\eventanchor{EQ-1662-ZHENG-DEATH-PRELUDE}{「郑成功在台湾去世，郑氏集团…}永历十六年／康熙元年五月（1662年），「郑成功在台湾去世，郑氏集团进入继承与军政重组阶段」之前的条件、信息和争议已通过中央与地方材料显现，构成该事件可追踪的前史节点。核心取证：《清圣祖实录》康熙元年海防军报；《清史稿·郑成功传》；《巴达维亚城日记》1662年台湾消息与郑氏文书辑录；同时期地方档案、地方志、日记或对方材料应按具体地区补检。该项锚定的是前史条件与信息背景，不把条件直接写成唯一因果。来源保留：官修编年和地方材料的记录密度与立场并不相同；本行不将条件、传闻、呈报或后见解释直接等同为确定因果。
\eventanchor{EQ-1662-ZHENG-DEATH-DECISION}{围绕「郑成功在台湾去世，郑氏…}永历十六年／康熙元年五月（1662年），围绕「郑成功在台湾去世，郑氏集团进入继承与军政重组阶段」的命令、调度、照会或呈报形成独立的中央—地方行动文书链。核心取证：《清圣祖实录》康熙元年海防军报；《清史稿·郑成功传》；《巴达维亚城日记》1662年台湾消息与郑氏文书辑录。该项锚定的是命令或决策环节，命令抵达和结果仍须另外核验。来源保留：奏折、上谕、部议、军机档或条约可证明行政动作和机构立场，不自动证明所有地区、群体或对手按其意图行动。
\eventanchor{EQ-1662-ZHENG-DEATH-AFTERMATH}{「郑成功在台湾去世，郑氏集团…}永历十六年／康熙元年五月（1662年），「郑成功在台湾去世，郑氏集团进入继承与军政重组阶段」引发的善后、执行或地方回应构成可由不同材料追踪的后续节点。核心取证：《清圣祖实录》康熙元年海防军报；《清史稿·郑成功传》；《巴达维亚城日记》1662年台湾消息与郑氏文书辑录；相关地区的档案、志书、契约、报刊或对方材料。该项锚定的是后续追踪问题，影响范围须按地区与群体分别判断。来源保留：战后、改革后或条约后的记忆、地方记录和官方总结常具有事后选择性；后续影响须按地点、群体和时间分层，不作整体化推断。
\eventanchor{EQ-1663-XIAMEN-JINMEN-PRELUDE}{「清军联合部分郑氏降将攻取厦…}康熙二年（1663年），「清军联合部分郑氏降将攻取厦门、金门，郑氏主力退守台湾」之前的条件、信息和争议已通过中央与地方材料显现，构成该事件可追踪的前史节点。核心取证：《清圣祖实录》康熙二年福建军报；《清史稿·施琅传》；《巴达维亚城日记》1663年中国沿海消息；同时期地方档案、地方志、日记或对方材料应按具体地区补检。该项锚定的是前史条件与信息背景，不把条件直接写成唯一因果。来源保留：官修编年和地方材料的记录密度与立场并不相同；本行不将条件、传闻、呈报或后见解释直接等同为确定因果。
\eventanchor{EQ-1663-XIAMEN-JINMEN-DECISION}{围绕「清军联合部分郑氏降将攻…}康熙二年（1663年），围绕「清军联合部分郑氏降将攻取厦门、金门，郑氏主力退守台湾」的命令、调度、照会或呈报形成独立的中央—地方行动文书链。核心取证：《清圣祖实录》康熙二年福建军报；《清史稿·施琅传》；《巴达维亚城日记》1663年中国沿海消息。该项锚定的是命令或决策环节，命令抵达和结果仍须另外核验。来源保留：奏折、上谕、部议、军机档或条约可证明行政动作和机构立场，不自动证明所有地区、群体或对手按其意图行动。
\eventanchor{EQ-1663-XIAMEN-JINMEN-AFTERMATH}{「清军联合部分郑氏降将攻取厦…}康熙二年（1663年），「清军联合部分郑氏降将攻取厦门、金门，郑氏主力退守台湾」引发的善后、执行或地方回应构成可由不同材料追踪的后续节点。核心取证：《清圣祖实录》康熙二年福建军报；《清史稿·施琅传》；《巴达维亚城日记》1663年中国沿海消息；相关地区的档案、志书、契约、报刊或对方材料。该项锚定的是后续追踪问题，影响范围须按地区与群体分别判断。来源保留：战后、改革后或条约后的记忆、地方记录和官方总结常具有事后选择性；后续影响须按地点、群体和时间分层，不作整体化推断。
\eventanchor{EQ-1664-KUIZHOU-PRELUDE}{「清军攻破夔东据点，李来亨等…}康熙三年（1664年），「清军攻破夔东据点，李来亨等抗清武装失败，长江上游的明遗民军事中心瓦解」之前的条件、信息和争议已通过中央与地方材料显现，构成该事件可追踪的前史节点。核心取证：《清圣祖实录》康熙三年湖广四川军报；《清史稿·李来亨传》；《夔州府志》与南明遗民文集；同时期地方档案、地方志、日记或对方材料应按具体地区补检。该项锚定的是前史条件与信息背景，不把条件直接写成唯一因果。来源保留：官修编年和地方材料的记录密度与立场并不相同；本行不将条件、传闻、呈报或后见解释直接等同为确定因果。
\eventanchor{EQ-1664-KUIZHOU-DECISION}{围绕「清军攻破夔东据点，李来…}康熙三年（1664年），围绕「清军攻破夔东据点，李来亨等抗清武装失败，长江上游的明遗民军事中心瓦解」的命令、调度、照会或呈报形成独立的中央—地方行动文书链。核心取证：《清圣祖实录》康熙三年湖广四川军报；《清史稿·李来亨传》；《夔州府志》与南明遗民文集。该项锚定的是命令或决策环节，命令抵达和结果仍须另外核验。来源保留：奏折、上谕、部议、军机档或条约可证明行政动作和机构立场，不自动证明所有地区、群体或对手按其意图行动。
\eventanchor{EQ-1664-KUIZHOU-AFTERMATH}{「清军攻破夔东据点，李来亨等…}康熙三年（1664年），「清军攻破夔东据点，李来亨等抗清武装失败，长江上游的明遗民军事中心瓦解」引发的善后、执行或地方回应构成可由不同材料追踪的后续节点。核心取证：《清圣祖实录》康熙三年湖广四川军报；《清史稿·李来亨传》；《夔州府志》与南明遗民文集；相关地区的档案、志书、契约、报刊或对方材料。该项锚定的是后续追踪问题，影响范围须按地区与群体分别判断。来源保留：战后、改革后或条约后的记忆、地方记录和官方总结常具有事后选择性；后续影响须按地点、群体和时间分层，不作整体化推断。
\eventanchor{EQ-1665-REGENTS-CONSOLIDATION-PRELUDE}{「四辅政大臣以议政和旗务网络…}康熙四年（1665年），「四辅政大臣以议政和旗务网络处理军政事务，鳌拜在政务中的影响扩大」之前的条件、信息和争议已通过中央与地方材料显现，构成该事件可追踪的前史节点。核心取证：《清圣祖实录》康熙四年政务条；《清史稿·鳌拜传》；《康熙起居注》及内阁大库满汉题本；同时期地方档案、地方志、日记或对方材料应按具体地区补检。该项锚定的是前史条件与信息背景，不把条件直接写成唯一因果。来源保留：官修编年和地方材料的记录密度与立场并不相同；本行不将条件、传闻、呈报或后见解释直接等同为确定因果。
\eventanchor{EQ-1665-REGENTS-CONSOLIDATION-DECISION}{围绕「四辅政大臣以议政和旗务…}康熙四年（1665年），围绕「四辅政大臣以议政和旗务网络处理军政事务，鳌拜在政务中的影响扩大」的命令、调度、照会或呈报形成独立的中央—地方行动文书链。核心取证：《清圣祖实录》康熙四年政务条；《清史稿·鳌拜传》；《康熙起居注》及内阁大库满汉题本。该项锚定的是命令或决策环节，命令抵达和结果仍须另外核验。来源保留：奏折、上谕、部议、军机档或条约可证明行政动作和机构立场，不自动证明所有地区、群体或对手按其意图行动。
\eventanchor{EQ-1665-REGENTS-CONSOLIDATION-AFTERMATH}{「四辅政大臣以议政和旗务网络…}康熙四年（1665年），「四辅政大臣以议政和旗务网络处理军政事务，鳌拜在政务中的影响扩大」引发的善后、执行或地方回应构成可由不同材料追踪的后续节点。核心取证：《清圣祖实录》康熙四年政务条；《清史稿·鳌拜传》；《康熙起居注》及内阁大库满汉题本；相关地区的档案、志书、契约、报刊或对方材料。该项锚定的是后续追踪问题，影响范围须按地区与群体分别判断。来源保留：战后、改革后或条约后的记忆、地方记录和官方总结常具有事后选择性；后续影响须按地点、群体和时间分层，不作整体化推断。
\eventanchor{EQ-1667-SONI-DEATH-PRELUDE}{「首辅索尼去世，四辅政大臣的…}康熙六年（1667年），「首辅索尼去世，四辅政大臣的均衡被打破」之前的条件、信息和争议已通过中央与地方材料显现，构成该事件可追踪的前史节点。核心取证：《清圣祖实录》康熙六年；《清史稿·索尼传》；《康熙起居注》康熙六年丧仪与议政记录；同时期地方档案、地方志、日记或对方材料应按具体地区补检。该项锚定的是前史条件与信息背景，不把条件直接写成唯一因果。来源保留：官修编年和地方材料的记录密度与立场并不相同；本行不将条件、传闻、呈报或后见解释直接等同为确定因果。
\eventanchor{EQ-1667-SONI-DEATH-DECISION}{围绕「首辅索尼去世，四辅政大…}康熙六年（1667年），围绕「首辅索尼去世，四辅政大臣的均衡被打破」的命令、调度、照会或呈报形成独立的中央—地方行动文书链。核心取证：《清圣祖实录》康熙六年；《清史稿·索尼传》；《康熙起居注》康熙六年丧仪与议政记录。该项锚定的是命令或决策环节，命令抵达和结果仍须另外核验。来源保留：奏折、上谕、部议、军机档或条约可证明行政动作和机构立场，不自动证明所有地区、群体或对手按其意图行动。
\eventanchor{EQ-1667-SONI-DEATH-AFTERMATH}{「首辅索尼去世，四辅政大臣的…}康熙六年（1667年），「首辅索尼去世，四辅政大臣的均衡被打破」引发的善后、执行或地方回应构成可由不同材料追踪的后续节点。核心取证：《清圣祖实录》康熙六年；《清史稿·索尼传》；《康熙起居注》康熙六年丧仪与议政记录；相关地区的档案、志书、契约、报刊或对方材料。该项锚定的是后续追踪问题，影响范围须按地区与群体分别判断。来源保留：战后、改革后或条约后的记忆、地方记录和官方总结常具有事后选择性；后续影响须按地点、群体和时间分层，不作整体化推断。
\eventanchor{EQ-1669-OBOI-ARREST-PRELUDE}{「康熙帝拘捕鳌拜并审理其党羽…}康熙八年五月（1669年），「康熙帝拘捕鳌拜并审理其党羽，结束辅政大臣实际主政的局面」之前的条件、信息和争议已通过中央与地方材料显现，构成该事件可追踪的前史节点。核心取证：《清圣祖实录》康熙八年五月；《清史稿·鳌拜传》；《康熙起居注》与鳌拜案内阁档；同时期地方档案、地方志、日记或对方材料应按具体地区补检。该项锚定的是前史条件与信息背景，不把条件直接写成唯一因果。来源保留：官修编年和地方材料的记录密度与立场并不相同；本行不将条件、传闻、呈报或后见解释直接等同为确定因果。
\eventanchor{EQ-1669-OBOI-ARREST-DECISION}{围绕「康熙帝拘捕鳌拜并审理其…}康熙八年五月（1669年），围绕「康熙帝拘捕鳌拜并审理其党羽，结束辅政大臣实际主政的局面」的命令、调度、照会或呈报形成独立的中央—地方行动文书链。核心取证：《清圣祖实录》康熙八年五月；《清史稿·鳌拜传》；《康熙起居注》与鳌拜案内阁档。该项锚定的是命令或决策环节，命令抵达和结果仍须另外核验。来源保留：奏折、上谕、部议、军机档或条约可证明行政动作和机构立场，不自动证明所有地区、群体或对手按其意图行动。
\eventanchor{EQ-1669-OBOI-ARREST-AFTERMATH}{「康熙帝拘捕鳌拜并审理其党羽…}康熙八年五月（1669年），「康熙帝拘捕鳌拜并审理其党羽，结束辅政大臣实际主政的局面」引发的善后、执行或地方回应构成可由不同材料追踪的后续节点。核心取证：《清圣祖实录》康熙八年五月；《清史稿·鳌拜传》；《康熙起居注》与鳌拜案内阁档；相关地区的档案、志书、契约、报刊或对方材料。该项锚定的是后续追踪问题，影响范围须按地区与群体分别判断。来源保留：战后、改革后或条约后的记忆、地方记录和官方总结常具有事后选择性；后续影响须按地点、群体和时间分层，不作整体化推断。
\subsection{1670—1679}
\eventanchor{EQ-1670-SACRED-EDICT-PRELUDE}{「康熙帝颁布十六条圣谕，规定…}康熙九年（1670年），「康熙帝颁布十六条圣谕，规定由地方官和民众遵循的伦理、赋役与秩序训导框架」之前的条件、信息和争议已通过中央与地方材料显现，构成该事件可追踪的前史节点。核心取证：《清圣祖实录》康熙九年十月；《清史稿·圣祖本纪》；《大清会典事例》圣谕条与各省地方志宣讲记录；同时期地方档案、地方志、日记或对方材料应按具体地区补检。该项锚定的是前史条件与信息背景，不把条件直接写成唯一因果。来源保留：官修编年和地方材料的记录密度与立场并不相同；本行不将条件、传闻、呈报或后见解释直接等同为确定因果。
\eventanchor{EQ-1670-SACRED-EDICT-DECISION}{围绕「康熙帝颁布十六条圣谕，…}康熙九年（1670年），围绕「康熙帝颁布十六条圣谕，规定由地方官和民众遵循的伦理、赋役与秩序训导框架」的命令、调度、照会或呈报形成独立的中央—地方行动文书链。核心取证：《清圣祖实录》康熙九年十月；《清史稿·圣祖本纪》；《大清会典事例》圣谕条与各省地方志宣讲记录。该项锚定的是命令或决策环节，命令抵达和结果仍须另外核验。来源保留：奏折、上谕、部议、军机档或条约可证明行政动作和机构立场，不自动证明所有地区、群体或对手按其意图行动。
\eventanchor{EQ-1670-SACRED-EDICT-AFTERMATH}{「康熙帝颁布十六条圣谕，规定…}康熙九年（1670年），「康熙帝颁布十六条圣谕，规定由地方官和民众遵循的伦理、赋役与秩序训导框架」引发的善后、执行或地方回应构成可由不同材料追踪的后续节点。核心取证：《清圣祖实录》康熙九年十月；《清史稿·圣祖本纪》；《大清会典事例》圣谕条与各省地方志宣讲记录；相关地区的档案、志书、契约、报刊或对方材料。该项锚定的是后续追踪问题，影响范围须按地区与群体分别判断。来源保留：战后、改革后或条约后的记忆、地方记录和官方总结常具有事后选择性；后续影响须按地点、群体和时间分层，不作整体化推断。
\eventanchor{EQ-1673-FEUDATORY-ABOLITION-PRELUDE}{「康熙帝批准撤除吴三桂、尚可…}康熙十二年三月（1673年），「康熙帝批准撤除吴三桂、尚可喜、耿继茂三藩的请求或安排，触发藩镇利益危机」之前的条件、信息和争议已通过中央与地方材料显现，构成该事件可追踪的前史节点。核心取证：《清圣祖实录》康熙十二年三月；《清史稿·三藩传》；内阁大库康熙十二年撤藩题本与《平定三逆方略》；同时期地方档案、地方志、日记或对方材料应按具体地区补检。该项锚定的是前史条件与信息背景，不把条件直接写成唯一因果。来源保留：官修编年和地方材料的记录密度与立场并不相同；本行不将条件、传闻、呈报或后见解释直接等同为确定因果。
\eventanchor{EQ-1673-FEUDATORY-ABOLITION-DECISION}{围绕「康熙帝批准撤除吴三桂、…}康熙十二年三月（1673年），围绕「康熙帝批准撤除吴三桂、尚可喜、耿继茂三藩的请求或安排，触发藩镇利益危机」的命令、调度、照会或呈报形成独立的中央—地方行动文书链。核心取证：《清圣祖实录》康熙十二年三月；《清史稿·三藩传》；内阁大库康熙十二年撤藩题本与《平定三逆方略》。该项锚定的是命令或决策环节，命令抵达和结果仍须另外核验。来源保留：奏折、上谕、部议、军机档或条约可证明行政动作和机构立场，不自动证明所有地区、群体或对手按其意图行动。
\eventanchor{EQ-1673-FEUDATORY-ABOLITION-AFTERMATH}{「康熙帝批准撤除吴三桂、尚可…}康熙十二年三月（1673年），「康熙帝批准撤除吴三桂、尚可喜、耿继茂三藩的请求或安排，触发藩镇利益危机」引发的善后、执行或地方回应构成可由不同材料追踪的后续节点。核心取证：《清圣祖实录》康熙十二年三月；《清史稿·三藩传》；内阁大库康熙十二年撤藩题本与《平定三逆方略》；相关地区的档案、志书、契约、报刊或对方材料。该项锚定的是后续追踪问题，影响范围须按地区与群体分别判断。来源保留：战后、改革后或条约后的记忆、地方记录和官方总结常具有事后选择性；后续影响须按地点、群体和时间分层，不作整体化推断。
\eventanchor{EQ-1673-WU-REVOLT-PRELUDE}{「吴三桂在云南起兵反清，三藩…}康熙十二年十一月（1673年），「吴三桂在云南起兵反清，三藩之乱正式爆发」之前的条件、信息和争议已通过中央与地方材料显现，构成该事件可追踪的前史节点。核心取证：《清圣祖实录》康熙十二年十一月；《清史稿·吴三桂传》；《平定三逆方略》卷首及云南地方志；同时期地方档案、地方志、日记或对方材料应按具体地区补检。该项锚定的是前史条件与信息背景，不把条件直接写成唯一因果。来源保留：官修编年和地方材料的记录密度与立场并不相同；本行不将条件、传闻、呈报或后见解释直接等同为确定因果。
\eventanchor{EQ-1673-WU-REVOLT-DECISION}{围绕「吴三桂在云南起兵反清，…}康熙十二年十一月（1673年），围绕「吴三桂在云南起兵反清，三藩之乱正式爆发」的命令、调度、照会或呈报形成独立的中央—地方行动文书链。核心取证：《清圣祖实录》康熙十二年十一月；《清史稿·吴三桂传》；《平定三逆方略》卷首及云南地方志。该项锚定的是命令或决策环节，命令抵达和结果仍须另外核验。来源保留：奏折、上谕、部议、军机档或条约可证明行政动作和机构立场，不自动证明所有地区、群体或对手按其意图行动。
\eventanchor{EQ-1673-WU-REVOLT-AFTERMATH}{「吴三桂在云南起兵反清，三藩…}康熙十二年十一月（1673年），「吴三桂在云南起兵反清，三藩之乱正式爆发」引发的善后、执行或地方回应构成可由不同材料追踪的后续节点。核心取证：《清圣祖实录》康熙十二年十一月；《清史稿·吴三桂传》；《平定三逆方略》卷首及云南地方志；相关地区的档案、志书、契约、报刊或对方材料。该项锚定的是后续追踪问题，影响范围须按地区与群体分别判断。来源保留：战后、改革后或条约后的记忆、地方记录和官方总结常具有事后选择性；后续影响须按地点、群体和时间分层，不作整体化推断。
\eventanchor{EQ-1674-GENG-JINGZHONG-PRELUDE}{「靖南王耿精忠在福建反清，东…}康熙十三年（1674年），「靖南王耿精忠在福建反清，东南海陆战场与郑氏关系重新组合」之前的条件、信息和争议已通过中央与地方材料显现，构成该事件可追踪的前史节点。核心取证：《清圣祖实录》康熙十三年福建军报；《清史稿·耿精忠传》；《平定三逆方略》与《福建通志》；同时期地方档案、地方志、日记或对方材料应按具体地区补检。该项锚定的是前史条件与信息背景，不把条件直接写成唯一因果。来源保留：官修编年和地方材料的记录密度与立场并不相同；本行不将条件、传闻、呈报或后见解释直接等同为确定因果。
\eventanchor{EQ-1674-GENG-JINGZHONG-DECISION}{围绕「靖南王耿精忠在福建反清…}康熙十三年（1674年），围绕「靖南王耿精忠在福建反清，东南海陆战场与郑氏关系重新组合」的命令、调度、照会或呈报形成独立的中央—地方行动文书链。核心取证：《清圣祖实录》康熙十三年福建军报；《清史稿·耿精忠传》；《平定三逆方略》与《福建通志》。该项锚定的是命令或决策环节，命令抵达和结果仍须另外核验。来源保留：奏折、上谕、部议、军机档或条约可证明行政动作和机构立场，不自动证明所有地区、群体或对手按其意图行动。
\eventanchor{EQ-1674-GENG-JINGZHONG-AFTERMATH}{「靖南王耿精忠在福建反清，东…}康熙十三年（1674年），「靖南王耿精忠在福建反清，东南海陆战场与郑氏关系重新组合」引发的善后、执行或地方回应构成可由不同材料追踪的后续节点。核心取证：《清圣祖实录》康熙十三年福建军报；《清史稿·耿精忠传》；《平定三逆方略》与《福建通志》；相关地区的档案、志书、契约、报刊或对方材料。该项锚定的是后续追踪问题，影响范围须按地区与群体分别判断。来源保留：战后、改革后或条约后的记忆、地方记录和官方总结常具有事后选择性；后续影响须按地点、群体和时间分层，不作整体化推断。
\eventanchor{EQ-1674-WANG-FUCHEN-PRELUDE}{「陕西提督王辅臣据平凉反清，…}康熙十三年（1674年），「陕西提督王辅臣据平凉反清，西北军政和通往甘肃的交通受威胁」之前的条件、信息和争议已通过中央与地方材料显现，构成该事件可追踪的前史节点。核心取证：《清圣祖实录》康熙十三年陕西军报；《清史稿·王辅臣传》；《平定三逆方略》与《平凉府志》；同时期地方档案、地方志、日记或对方材料应按具体地区补检。该项锚定的是前史条件与信息背景，不把条件直接写成唯一因果。来源保留：官修编年和地方材料的记录密度与立场并不相同；本行不将条件、传闻、呈报或后见解释直接等同为确定因果。
\eventanchor{EQ-1674-WANG-FUCHEN-DECISION}{围绕「陕西提督王辅臣据平凉反…}康熙十三年（1674年），围绕「陕西提督王辅臣据平凉反清，西北军政和通往甘肃的交通受威胁」的命令、调度、照会或呈报形成独立的中央—地方行动文书链。核心取证：《清圣祖实录》康熙十三年陕西军报；《清史稿·王辅臣传》；《平定三逆方略》与《平凉府志》。该项锚定的是命令或决策环节，命令抵达和结果仍须另外核验。来源保留：奏折、上谕、部议、军机档或条约可证明行政动作和机构立场，不自动证明所有地区、群体或对手按其意图行动。
\eventanchor{EQ-1674-WANG-FUCHEN-AFTERMATH}{「陕西提督王辅臣据平凉反清，…}康熙十三年（1674年），「陕西提督王辅臣据平凉反清，西北军政和通往甘肃的交通受威胁」引发的善后、执行或地方回应构成可由不同材料追踪的后续节点。核心取证：《清圣祖实录》康熙十三年陕西军报；《清史稿·王辅臣传》；《平定三逆方略》与《平凉府志》；相关地区的档案、志书、契约、报刊或对方材料。该项锚定的是后续追踪问题，影响范围须按地区与群体分别判断。来源保留：战后、改革后或条约后的记忆、地方记录和官方总结常具有事后选择性；后续影响须按地点、群体和时间分层，不作整体化推断。
\eventanchor{EQ-1674-ZHENG-JING-FUJIAN-PRELUDE}{「郑经自台湾出兵福建，取得漳…}康熙十三年（1674年），「郑经自台湾出兵福建，取得漳泉沿海若干据点并介入三藩战争」之前的条件、信息和争议已通过中央与地方材料显现，构成该事件可追踪的前史节点。核心取证：《清圣祖实录》康熙十三年福建海防军报；《清史稿·郑成功传附郑经》；《台湾外纪》与《福建通志》；同时期地方档案、地方志、日记或对方材料应按具体地区补检。该项锚定的是前史条件与信息背景，不把条件直接写成唯一因果。来源保留：官修编年和地方材料的记录密度与立场并不相同；本行不将条件、传闻、呈报或后见解释直接等同为确定因果。
\eventanchor{EQ-1674-ZHENG-JING-FUJIAN-DECISION}{围绕「郑经自台湾出兵福建，取…}康熙十三年（1674年），围绕「郑经自台湾出兵福建，取得漳泉沿海若干据点并介入三藩战争」的命令、调度、照会或呈报形成独立的中央—地方行动文书链。核心取证：《清圣祖实录》康熙十三年福建海防军报；《清史稿·郑成功传附郑经》；《台湾外纪》与《福建通志》。该项锚定的是命令或决策环节，命令抵达和结果仍须另外核验。来源保留：奏折、上谕、部议、军机档或条约可证明行政动作和机构立场，不自动证明所有地区、群体或对手按其意图行动。
\eventanchor{EQ-1674-ZHENG-JING-FUJIAN-AFTERMATH}{「郑经自台湾出兵福建，取得漳…}康熙十三年（1674年），「郑经自台湾出兵福建，取得漳泉沿海若干据点并介入三藩战争」引发的善后、执行或地方回应构成可由不同材料追踪的后续节点。核心取证：《清圣祖实录》康熙十三年福建海防军报；《清史稿·郑成功传附郑经》；《台湾外纪》与《福建通志》；相关地区的档案、志书、契约、报刊或对方材料。该项锚定的是后续追踪问题，影响范围须按地区与群体分别判断。来源保留：战后、改革后或条约后的记忆、地方记录和官方总结常具有事后选择性；后续影响须按地点、群体和时间分层，不作整体化推断。
\eventanchor{EQ-1675-SHANG-ZHIXIN-PRELUDE}{「平南王尚之信在广州控制军政…}康熙十四年（1675年），「平南王尚之信在广州控制军政并与反清势力周旋，两广局势动荡」之前的条件、信息和争议已通过中央与地方材料显现，构成该事件可追踪的前史节点。核心取证：《清圣祖实录》康熙十四年广东军报；《清史稿·尚之信传》；《平定三逆方略》与《广东通志》；同时期地方档案、地方志、日记或对方材料应按具体地区补检。该项锚定的是前史条件与信息背景，不把条件直接写成唯一因果。来源保留：官修编年和地方材料的记录密度与立场并不相同；本行不将条件、传闻、呈报或后见解释直接等同为确定因果。
\eventanchor{EQ-1675-SHANG-ZHIXIN-DECISION}{围绕「平南王尚之信在广州控制…}康熙十四年（1675年），围绕「平南王尚之信在广州控制军政并与反清势力周旋，两广局势动荡」的命令、调度、照会或呈报形成独立的中央—地方行动文书链。核心取证：《清圣祖实录》康熙十四年广东军报；《清史稿·尚之信传》；《平定三逆方略》与《广东通志》。该项锚定的是命令或决策环节，命令抵达和结果仍须另外核验。来源保留：奏折、上谕、部议、军机档或条约可证明行政动作和机构立场，不自动证明所有地区、群体或对手按其意图行动。
\eventanchor{EQ-1675-SHANG-ZHIXIN-AFTERMATH}{「平南王尚之信在广州控制军政…}康熙十四年（1675年），「平南王尚之信在广州控制军政并与反清势力周旋，两广局势动荡」引发的善后、执行或地方回应构成可由不同材料追踪的后续节点。核心取证：《清圣祖实录》康熙十四年广东军报；《清史稿·尚之信传》；《平定三逆方略》与《广东通志》；相关地区的档案、志书、契约、报刊或对方材料。该项锚定的是后续追踪问题，影响范围须按地区与群体分别判断。来源保留：战后、改革后或条约后的记忆、地方记录和官方总结常具有事后选择性；后续影响须按地点、群体和时间分层，不作整体化推断。
\eventanchor{EQ-1676-GENG-SURRENDER-PRELUDE}{「耿精忠向清廷投降，福建藩镇…}康熙十五年（1676年），「耿精忠向清廷投降，福建藩镇叛乱的核心力量瓦解」之前的条件、信息和争议已通过中央与地方材料显现，构成该事件可追踪的前史节点。核心取证：《清圣祖实录》康熙十五年；《清史稿·耿精忠传》；内阁大库康熙十五年福建招抚题本与《福建通志》；同时期地方档案、地方志、日记或对方材料应按具体地区补检。该项锚定的是前史条件与信息背景，不把条件直接写成唯一因果。来源保留：官修编年和地方材料的记录密度与立场并不相同；本行不将条件、传闻、呈报或后见解释直接等同为确定因果。
\eventanchor{EQ-1676-GENG-SURRENDER-DECISION}{围绕「耿精忠向清廷投降，福建…}康熙十五年（1676年），围绕「耿精忠向清廷投降，福建藩镇叛乱的核心力量瓦解」的命令、调度、照会或呈报形成独立的中央—地方行动文书链。核心取证：《清圣祖实录》康熙十五年；《清史稿·耿精忠传》；内阁大库康熙十五年福建招抚题本与《福建通志》。该项锚定的是命令或决策环节，命令抵达和结果仍须另外核验。来源保留：奏折、上谕、部议、军机档或条约可证明行政动作和机构立场，不自动证明所有地区、群体或对手按其意图行动。
\eventanchor{EQ-1676-GENG-SURRENDER-AFTERMATH}{「耿精忠向清廷投降，福建藩镇…}康熙十五年（1676年），「耿精忠向清廷投降，福建藩镇叛乱的核心力量瓦解」引发的善后、执行或地方回应构成可由不同材料追踪的后续节点。核心取证：《清圣祖实录》康熙十五年；《清史稿·耿精忠传》；内阁大库康熙十五年福建招抚题本与《福建通志》；相关地区的档案、志书、契约、报刊或对方材料。该项锚定的是后续追踪问题，影响范围须按地区与群体分别判断。来源保留：战后、改革后或条约后的记忆、地方记录和官方总结常具有事后选择性；后续影响须按地点、群体和时间分层，不作整体化推断。
\eventanchor{EQ-1676-WUCHANG-PRELUDE}{「清军收复武昌等长江中游要地…}康熙十五年（1676年），「清军收复武昌等长江中游要地，吴军向湖南、广西方向收缩」之前的条件、信息和争议已通过中央与地方材料显现，构成该事件可追踪的前史节点。核心取证：《清圣祖实录》康熙十五年湖广军报；《清史稿·赵良栋传》；《平定三逆方略》与《武昌府志》；同时期地方档案、地方志、日记或对方材料应按具体地区补检。该项锚定的是前史条件与信息背景，不把条件直接写成唯一因果。来源保留：官修编年和地方材料的记录密度与立场并不相同；本行不将条件、传闻、呈报或后见解释直接等同为确定因果。
\eventanchor{EQ-1676-WUCHANG-DECISION}{围绕「清军收复武昌等长江中游…}康熙十五年（1676年），围绕「清军收复武昌等长江中游要地，吴军向湖南、广西方向收缩」的命令、调度、照会或呈报形成独立的中央—地方行动文书链。核心取证：《清圣祖实录》康熙十五年湖广军报；《清史稿·赵良栋传》；《平定三逆方略》与《武昌府志》。该项锚定的是命令或决策环节，命令抵达和结果仍须另外核验。来源保留：奏折、上谕、部议、军机档或条约可证明行政动作和机构立场，不自动证明所有地区、群体或对手按其意图行动。
\eventanchor{EQ-1676-WUCHANG-AFTERMATH}{「清军收复武昌等长江中游要地…}康熙十五年（1676年），「清军收复武昌等长江中游要地，吴军向湖南、广西方向收缩」引发的善后、执行或地方回应构成可由不同材料追踪的后续节点。核心取证：《清圣祖实录》康熙十五年湖广军报；《清史稿·赵良栋传》；《平定三逆方略》与《武昌府志》；相关地区的档案、志书、契约、报刊或对方材料。该项锚定的是后续追踪问题，影响范围须按地区与群体分别判断。来源保留：战后、改革后或条约后的记忆、地方记录和官方总结常具有事后选择性；后续影响须按地点、群体和时间分层，不作整体化推断。
\eventanchor{EQ-1677-WU-ZHOU-PRELUDE}{「吴三桂在衡阳称帝，国号周，…}康熙十六年（1677年），「吴三桂在衡阳称帝，国号周，将撤藩战争转为公开争夺最高政治名义」之前的条件、信息和争议已通过中央与地方材料显现，构成该事件可追踪的前史节点。核心取证：《清圣祖实录》康熙十六年湖南军报；《清史稿·吴三桂传》；《平定三逆方略》及吴周檄文辑录；同时期地方档案、地方志、日记或对方材料应按具体地区补检。该项锚定的是前史条件与信息背景，不把条件直接写成唯一因果。来源保留：官修编年和地方材料的记录密度与立场并不相同；本行不将条件、传闻、呈报或后见解释直接等同为确定因果。
\eventanchor{EQ-1677-WU-ZHOU-DECISION}{围绕「吴三桂在衡阳称帝，国号…}康熙十六年（1677年），围绕「吴三桂在衡阳称帝，国号周，将撤藩战争转为公开争夺最高政治名义」的命令、调度、照会或呈报形成独立的中央—地方行动文书链。核心取证：《清圣祖实录》康熙十六年湖南军报；《清史稿·吴三桂传》；《平定三逆方略》及吴周檄文辑录。该项锚定的是命令或决策环节，命令抵达和结果仍须另外核验。来源保留：奏折、上谕、部议、军机档或条约可证明行政动作和机构立场，不自动证明所有地区、群体或对手按其意图行动。
\eventanchor{EQ-1677-WU-ZHOU-AFTERMATH}{「吴三桂在衡阳称帝，国号周，…}康熙十六年（1677年），「吴三桂在衡阳称帝，国号周，将撤藩战争转为公开争夺最高政治名义」引发的善后、执行或地方回应构成可由不同材料追踪的后续节点。核心取证：《清圣祖实录》康熙十六年湖南军报；《清史稿·吴三桂传》；《平定三逆方略》及吴周檄文辑录；相关地区的档案、志书、契约、报刊或对方材料。该项锚定的是后续追踪问题，影响范围须按地区与群体分别判断。来源保留：战后、改革后或条约后的记忆、地方记录和官方总结常具有事后选择性；后续影响须按地点、群体和时间分层，不作整体化推断。
\eventanchor{EQ-1678-WU-DEATH-PRELUDE}{「吴三桂去世，吴世璠继承吴周…}康熙十七年八月（1678年），「吴三桂去世，吴世璠继承吴周政权」之前的条件、信息和争议已通过中央与地方材料显现，构成该事件可追踪的前史节点。核心取证：《清圣祖实录》康熙十七年八月；《清史稿·吴三桂传》；《平定三逆方略》与湖南地方志；同时期地方档案、地方志、日记或对方材料应按具体地区补检。该项锚定的是前史条件与信息背景，不把条件直接写成唯一因果。来源保留：官修编年和地方材料的记录密度与立场并不相同；本行不将条件、传闻、呈报或后见解释直接等同为确定因果。
\eventanchor{EQ-1678-WU-DEATH-DECISION}{围绕「吴三桂去世，吴世璠继承…}康熙十七年八月（1678年），围绕「吴三桂去世，吴世璠继承吴周政权」的命令、调度、照会或呈报形成独立的中央—地方行动文书链。核心取证：《清圣祖实录》康熙十七年八月；《清史稿·吴三桂传》；《平定三逆方略》与湖南地方志。该项锚定的是命令或决策环节，命令抵达和结果仍须另外核验。来源保留：奏折、上谕、部议、军机档或条约可证明行政动作和机构立场，不自动证明所有地区、群体或对手按其意图行动。
\eventanchor{EQ-1678-WU-DEATH-AFTERMATH}{「吴三桂去世，吴世璠继承吴周…}康熙十七年八月（1678年），「吴三桂去世，吴世璠继承吴周政权」引发的善后、执行或地方回应构成可由不同材料追踪的后续节点。核心取证：《清圣祖实录》康熙十七年八月；《清史稿·吴三桂传》；《平定三逆方略》与湖南地方志；相关地区的档案、志书、契约、报刊或对方材料。该项锚定的是后续追踪问题，影响范围须按地区与群体分别判断。来源保留：战后、改革后或条约后的记忆、地方记录和官方总结常具有事后选择性；后续影响须按地点、群体和时间分层，不作整体化推断。
\eventanchor{EQ-1679-HUNAN-RECOVERY-PRELUDE}{「清军在湖南持续推进并夺回若…}康熙十八年（1679年），「清军在湖南持续推进并夺回若干府县，吴周军向贵州、广西收缩」之前的条件、信息和争议已通过中央与地方材料显现，构成该事件可追踪的前史节点。核心取证：《清圣祖实录》康熙十八年湖南军报；《清史稿·赵良栋传》；《平定三逆方略》及《湖南通志》；同时期地方档案、地方志、日记或对方材料应按具体地区补检。该项锚定的是前史条件与信息背景，不把条件直接写成唯一因果。来源保留：官修编年和地方材料的记录密度与立场并不相同；本行不将条件、传闻、呈报或后见解释直接等同为确定因果。
\eventanchor{EQ-1679-HUNAN-RECOVERY-DECISION}{围绕「清军在湖南持续推进并夺…}康熙十八年（1679年），围绕「清军在湖南持续推进并夺回若干府县，吴周军向贵州、广西收缩」的命令、调度、照会或呈报形成独立的中央—地方行动文书链。核心取证：《清圣祖实录》康熙十八年湖南军报；《清史稿·赵良栋传》；《平定三逆方略》及《湖南通志》。该项锚定的是命令或决策环节，命令抵达和结果仍须另外核验。来源保留：奏折、上谕、部议、军机档或条约可证明行政动作和机构立场，不自动证明所有地区、群体或对手按其意图行动。
\eventanchor{EQ-1679-HUNAN-RECOVERY-AFTERMATH}{「清军在湖南持续推进并夺回若…}康熙十八年（1679年），「清军在湖南持续推进并夺回若干府县，吴周军向贵州、广西收缩」引发的善后、执行或地方回应构成可由不同材料追踪的后续节点。核心取证：《清圣祖实录》康熙十八年湖南军报；《清史稿·赵良栋传》；《平定三逆方略》及《湖南通志》；相关地区的档案、志书、契约、报刊或对方材料。该项锚定的是后续追踪问题，影响范围须按地区与群体分别判断。来源保留：战后、改革后或条约后的记忆、地方记录和官方总结常具有事后选择性；后续影响须按地点、群体和时间分层，不作整体化推断。
\subsection{1680—1689}
\eventanchor{EQ-1680-SICHUAN-RECOVERY-PRELUDE}{「清军清理四川通往云南贵州的…}康熙十九年（1680年），「清军清理四川通往云南贵州的要道并恢复主要府城控制」之前的条件、信息和争议已通过中央与地方材料显现，构成该事件可追踪的前史节点。核心取证：《清圣祖实录》康熙十九年四川军报；《清史稿·傅弘烈传》；《四川通志》与内阁大库四川粮饷题本；同时期地方档案、地方志、日记或对方材料应按具体地区补检。该项锚定的是前史条件与信息背景，不把条件直接写成唯一因果。来源保留：官修编年和地方材料的记录密度与立场并不相同；本行不将条件、传闻、呈报或后见解释直接等同为确定因果。
\eventanchor{EQ-1680-SICHUAN-RECOVERY-DECISION}{围绕「清军清理四川通往云南贵…}康熙十九年（1680年），围绕「清军清理四川通往云南贵州的要道并恢复主要府城控制」的命令、调度、照会或呈报形成独立的中央—地方行动文书链。核心取证：《清圣祖实录》康熙十九年四川军报；《清史稿·傅弘烈传》；《四川通志》与内阁大库四川粮饷题本。该项锚定的是命令或决策环节，命令抵达和结果仍须另外核验。来源保留：奏折、上谕、部议、军机档或条约可证明行政动作和机构立场，不自动证明所有地区、群体或对手按其意图行动。
\eventanchor{EQ-1680-SICHUAN-RECOVERY-AFTERMATH}{「清军清理四川通往云南贵州的…}康熙十九年（1680年），「清军清理四川通往云南贵州的要道并恢复主要府城控制」引发的善后、执行或地方回应构成可由不同材料追踪的后续节点。核心取证：《清圣祖实录》康熙十九年四川军报；《清史稿·傅弘烈传》；《四川通志》与内阁大库四川粮饷题本；相关地区的档案、志书、契约、报刊或对方材料。该项锚定的是后续追踪问题，影响范围须按地区与群体分别判断。来源保留：战后、改革后或条约后的记忆、地方记录和官方总结常具有事后选择性；后续影响须按地点、群体和时间分层，不作整体化推断。
\eventanchor{EQ-1681-KUNMING-PRELUDE}{「清军攻入昆明，吴世璠自尽，…}康熙二十年十月（1681年），「清军攻入昆明，吴世璠自尽，三藩之乱的吴周政权终结」之前的条件、信息和争议已通过中央与地方材料显现，构成该事件可追踪的前史节点。核心取证：《清圣祖实录》康熙二十年十月；《清史稿·吴世璠传》；《平定三逆方略》与《云南通志》；同时期地方档案、地方志、日记或对方材料应按具体地区补检。该项锚定的是前史条件与信息背景，不把条件直接写成唯一因果。来源保留：官修编年和地方材料的记录密度与立场并不相同；本行不将条件、传闻、呈报或后见解释直接等同为确定因果。
\eventanchor{EQ-1681-KUNMING-DECISION}{围绕「清军攻入昆明，吴世璠自…}康熙二十年十月（1681年），围绕「清军攻入昆明，吴世璠自尽，三藩之乱的吴周政权终结」的命令、调度、照会或呈报形成独立的中央—地方行动文书链。核心取证：《清圣祖实录》康熙二十年十月；《清史稿·吴世璠传》；《平定三逆方略》与《云南通志》。该项锚定的是命令或决策环节，命令抵达和结果仍须另外核验。来源保留：奏折、上谕、部议、军机档或条约可证明行政动作和机构立场，不自动证明所有地区、群体或对手按其意图行动。
\eventanchor{EQ-1681-KUNMING-AFTERMATH}{「清军攻入昆明，吴世璠自尽，…}康熙二十年十月（1681年），「清军攻入昆明，吴世璠自尽，三藩之乱的吴周政权终结」引发的善后、执行或地方回应构成可由不同材料追踪的后续节点。核心取证：《清圣祖实录》康熙二十年十月；《清史稿·吴世璠传》；《平定三逆方略》与《云南通志》；相关地区的档案、志书、契约、报刊或对方材料。该项锚定的是后续追踪问题，影响范围须按地区与群体分别判断。来源保留：战后、改革后或条约后的记忆、地方记录和官方总结常具有事后选择性；后续影响须按地点、群体和时间分层，不作整体化推断。
\eventanchor{EQ-1681-ZHENG-JING-DEATH-PRELUDE}{「郑经去世，东宁发生继承冲突…}康熙二十年（1681年），「郑经去世，东宁发生继承冲突，郑克塽最终掌权」之前的条件、信息和争议已通过中央与地方材料显现，构成该事件可追踪的前史节点。核心取证：《清圣祖实录》康熙二十年台湾情报；《清史稿·郑成功传附郑经》；《台湾外纪》与荷兰商馆残存通信；同时期地方档案、地方志、日记或对方材料应按具体地区补检。该项锚定的是前史条件与信息背景，不把条件直接写成唯一因果。来源保留：官修编年和地方材料的记录密度与立场并不相同；本行不将条件、传闻、呈报或后见解释直接等同为确定因果。
\eventanchor{EQ-1681-ZHENG-JING-DEATH-DECISION}{围绕「郑经去世，东宁发生继承…}康熙二十年（1681年），围绕「郑经去世，东宁发生继承冲突，郑克塽最终掌权」的命令、调度、照会或呈报形成独立的中央—地方行动文书链。核心取证：《清圣祖实录》康熙二十年台湾情报；《清史稿·郑成功传附郑经》；《台湾外纪》与荷兰商馆残存通信。该项锚定的是命令或决策环节，命令抵达和结果仍须另外核验。来源保留：奏折、上谕、部议、军机档或条约可证明行政动作和机构立场，不自动证明所有地区、群体或对手按其意图行动。
\eventanchor{EQ-1681-ZHENG-JING-DEATH-AFTERMATH}{「郑经去世，东宁发生继承冲突…}康熙二十年（1681年），「郑经去世，东宁发生继承冲突，郑克塽最终掌权」引发的善后、执行或地方回应构成可由不同材料追踪的后续节点。核心取证：《清圣祖实录》康熙二十年台湾情报；《清史稿·郑成功传附郑经》；《台湾外纪》与荷兰商馆残存通信；相关地区的档案、志书、契约、报刊或对方材料。该项锚定的是后续追踪问题，影响范围须按地区与群体分别判断。来源保留：战后、改革后或条约后的记忆、地方记录和官方总结常具有事后选择性；后续影响须按地点、群体和时间分层，不作整体化推断。
\eventanchor{EQ-1682-SHI-LANG-PRELUDE}{「清廷任命施琅统率福建水师，…}康熙二十一年（1682年），「清廷任命施琅统率福建水师，筹备进攻台湾」之前的条件、信息和争议已通过中央与地方材料显现，构成该事件可追踪的前史节点。核心取证：《清圣祖实录》康熙二十一年福建海防条；《清史稿·施琅传》；《康熙朝满文朱批奏折全译》康熙二十一年台湾军务奏折；同时期地方档案、地方志、日记或对方材料应按具体地区补检。该项锚定的是前史条件与信息背景，不把条件直接写成唯一因果。来源保留：官修编年和地方材料的记录密度与立场并不相同；本行不将条件、传闻、呈报或后见解释直接等同为确定因果。
\eventanchor{EQ-1682-SHI-LANG-DECISION}{围绕「清廷任命施琅统率福建水…}康熙二十一年（1682年），围绕「清廷任命施琅统率福建水师，筹备进攻台湾」的命令、调度、照会或呈报形成独立的中央—地方行动文书链。核心取证：《清圣祖实录》康熙二十一年福建海防条；《清史稿·施琅传》；《康熙朝满文朱批奏折全译》康熙二十一年台湾军务奏折。该项锚定的是命令或决策环节，命令抵达和结果仍须另外核验。来源保留：奏折、上谕、部议、军机档或条约可证明行政动作和机构立场，不自动证明所有地区、群体或对手按其意图行动。
\eventanchor{EQ-1682-SHI-LANG-AFTERMATH}{「清廷任命施琅统率福建水师，…}康熙二十一年（1682年），「清廷任命施琅统率福建水师，筹备进攻台湾」引发的善后、执行或地方回应构成可由不同材料追踪的后续节点。核心取证：《清圣祖实录》康熙二十一年福建海防条；《清史稿·施琅传》；《康熙朝满文朱批奏折全译》康熙二十一年台湾军务奏折；相关地区的档案、志书、契约、报刊或对方材料。该项锚定的是后续追踪问题，影响范围须按地区与群体分别判断。来源保留：战后、改革后或条约后的记忆、地方记录和官方总结常具有事后选择性；后续影响须按地点、群体和时间分层，不作整体化推断。
\eventanchor{EQ-1683-PENGHU-PRELUDE}{「施琅水师在澎湖击败郑氏舰队…}康熙二十二年六月（1683年），「施琅水师在澎湖击败郑氏舰队，取得通往台湾本岛的制海优势」之前的条件、信息和争议已通过中央与地方材料显现，构成该事件可追踪的前史节点。核心取证：《清圣祖实录》康熙二十二年六月；《清史稿·施琅传》；施琅《飞报澎湖大捷疏》与台湾地方志；同时期地方档案、地方志、日记或对方材料应按具体地区补检。该项锚定的是前史条件与信息背景，不把条件直接写成唯一因果。来源保留：官修编年和地方材料的记录密度与立场并不相同；本行不将条件、传闻、呈报或后见解释直接等同为确定因果。
\eventanchor{EQ-1683-PENGHU-DECISION}{围绕「施琅水师在澎湖击败郑氏…}康熙二十二年六月（1683年），围绕「施琅水师在澎湖击败郑氏舰队，取得通往台湾本岛的制海优势」的命令、调度、照会或呈报形成独立的中央—地方行动文书链。核心取证：《清圣祖实录》康熙二十二年六月；《清史稿·施琅传》；施琅《飞报澎湖大捷疏》与台湾地方志。该项锚定的是命令或决策环节，命令抵达和结果仍须另外核验。来源保留：奏折、上谕、部议、军机档或条约可证明行政动作和机构立场，不自动证明所有地区、群体或对手按其意图行动。
\eventanchor{EQ-1683-PENGHU-AFTERMATH}{「施琅水师在澎湖击败郑氏舰队…}康熙二十二年六月（1683年），「施琅水师在澎湖击败郑氏舰队，取得通往台湾本岛的制海优势」引发的善后、执行或地方回应构成可由不同材料追踪的后续节点。核心取证：《清圣祖实录》康熙二十二年六月；《清史稿·施琅传》；施琅《飞报澎湖大捷疏》与台湾地方志；相关地区的档案、志书、契约、报刊或对方材料。该项锚定的是后续追踪问题，影响范围须按地区与群体分别判断。来源保留：战后、改革后或条约后的记忆、地方记录和官方总结常具有事后选择性；后续影响须按地点、群体和时间分层，不作整体化推断。
\eventanchor{EQ-1683-TAIWAN-SURRENDER-PRELUDE}{「郑克塽向清廷投降，东宁政权…}康熙二十二年八月（1683年），「郑克塽向清廷投降，东宁政权结束」之前的条件、信息和争议已通过中央与地方材料显现，构成该事件可追踪的前史节点。核心取证：《清圣祖实录》康熙二十二年八月；《清史稿·郑成功传附郑克塽》；施琅奏折与《台湾府志》；同时期地方档案、地方志、日记或对方材料应按具体地区补检。该项锚定的是前史条件与信息背景，不把条件直接写成唯一因果。来源保留：官修编年和地方材料的记录密度与立场并不相同；本行不将条件、传闻、呈报或后见解释直接等同为确定因果。
\eventanchor{EQ-1683-TAIWAN-SURRENDER-DECISION}{围绕「郑克塽向清廷投降，东宁…}康熙二十二年八月（1683年），围绕「郑克塽向清廷投降，东宁政权结束」的命令、调度、照会或呈报形成独立的中央—地方行动文书链。核心取证：《清圣祖实录》康熙二十二年八月；《清史稿·郑成功传附郑克塽》；施琅奏折与《台湾府志》。该项锚定的是命令或决策环节，命令抵达和结果仍须另外核验。来源保留：奏折、上谕、部议、军机档或条约可证明行政动作和机构立场，不自动证明所有地区、群体或对手按其意图行动。
\eventanchor{EQ-1683-TAIWAN-SURRENDER-AFTERMATH}{「郑克塽向清廷投降，东宁政权…}康熙二十二年八月（1683年），「郑克塽向清廷投降，东宁政权结束」引发的善后、执行或地方回应构成可由不同材料追踪的后续节点。核心取证：《清圣祖实录》康熙二十二年八月；《清史稿·郑成功传附郑克塽》；施琅奏折与《台湾府志》；相关地区的档案、志书、契约、报刊或对方材料。该项锚定的是后续追踪问题，影响范围须按地区与群体分别判断。来源保留：战后、改革后或条约后的记忆、地方记录和官方总结常具有事后选择性；后续影响须按地点、群体和时间分层，不作整体化推断。
\eventanchor{EQ-1684-TAIWAN-PREFECTURE-PRELUDE}{「清廷设置台湾府及诸县，隶属…}康熙二十三年（1684年），「清廷设置台湾府及诸县，隶属福建省，开始以府县和驻防制度经营台湾」之前的条件、信息和争议已通过中央与地方材料显现，构成该事件可追踪的前史节点。核心取证：《清圣祖实录》康熙二十三年台湾设治条；《清史稿·地理志》；《台湾府志》与福建省档案相关奏议；同时期地方档案、地方志、日记或对方材料应按具体地区补检。该项锚定的是前史条件与信息背景，不把条件直接写成唯一因果。来源保留：官修编年和地方材料的记录密度与立场并不相同；本行不将条件、传闻、呈报或后见解释直接等同为确定因果。
\eventanchor{EQ-1684-TAIWAN-PREFECTURE-DECISION}{围绕「清廷设置台湾府及诸县，…}康熙二十三年（1684年），围绕「清廷设置台湾府及诸县，隶属福建省，开始以府县和驻防制度经营台湾」的命令、调度、照会或呈报形成独立的中央—地方行动文书链。核心取证：《清圣祖实录》康熙二十三年台湾设治条；《清史稿·地理志》；《台湾府志》与福建省档案相关奏议。该项锚定的是命令或决策环节，命令抵达和结果仍须另外核验。来源保留：奏折、上谕、部议、军机档或条约可证明行政动作和机构立场，不自动证明所有地区、群体或对手按其意图行动。
\eventanchor{EQ-1684-TAIWAN-PREFECTURE-AFTERMATH}{「清廷设置台湾府及诸县，隶属…}康熙二十三年（1684年），「清廷设置台湾府及诸县，隶属福建省，开始以府县和驻防制度经营台湾」引发的善后、执行或地方回应构成可由不同材料追踪的后续节点。核心取证：《清圣祖实录》康熙二十三年台湾设治条；《清史稿·地理志》；《台湾府志》与福建省档案相关奏议；相关地区的档案、志书、契约、报刊或对方材料。该项锚定的是后续追踪问题，影响范围须按地区与群体分别判断。来源保留：战后、改革后或条约后的记忆、地方记录和官方总结常具有事后选择性；后续影响须按地点、群体和时间分层，不作整体化推断。
\eventanchor{EQ-1684-SEA-BAN-LIFT-PRELUDE}{「清廷调整迁界与海禁政策，允…}康熙二十三年（1684年），「清廷调整迁界与海禁政策，允许沿海居民复界并逐步恢复民间海贸」之前的条件、信息和争议已通过中央与地方材料显现，构成该事件可追踪的前史节点。核心取证：《清圣祖实录》康熙二十三年复界海禁条；《清史稿·食货志》；《福建通志》《广东通志》沿海复界条；同时期地方档案、地方志、日记或对方材料应按具体地区补检。该项锚定的是前史条件与信息背景，不把条件直接写成唯一因果。来源保留：官修编年和地方材料的记录密度与立场并不相同；本行不将条件、传闻、呈报或后见解释直接等同为确定因果。
\eventanchor{EQ-1684-SEA-BAN-LIFT-DECISION}{围绕「清廷调整迁界与海禁政策…}康熙二十三年（1684年），围绕「清廷调整迁界与海禁政策，允许沿海居民复界并逐步恢复民间海贸」的命令、调度、照会或呈报形成独立的中央—地方行动文书链。核心取证：《清圣祖实录》康熙二十三年复界海禁条；《清史稿·食货志》；《福建通志》《广东通志》沿海复界条。该项锚定的是命令或决策环节，命令抵达和结果仍须另外核验。来源保留：奏折、上谕、部议、军机档或条约可证明行政动作和机构立场，不自动证明所有地区、群体或对手按其意图行动。
\eventanchor{EQ-1684-SEA-BAN-LIFT-AFTERMATH}{「清廷调整迁界与海禁政策，允…}康熙二十三年（1684年），「清廷调整迁界与海禁政策，允许沿海居民复界并逐步恢复民间海贸」引发的善后、执行或地方回应构成可由不同材料追踪的后续节点。核心取证：《清圣祖实录》康熙二十三年复界海禁条；《清史稿·食货志》；《福建通志》《广东通志》沿海复界条；相关地区的档案、志书、契约、报刊或对方材料。该项锚定的是后续追踪问题，影响范围须按地区与群体分别判断。来源保留：战后、改革后或条约后的记忆、地方记录和官方总结常具有事后选择性；后续影响须按地点、群体和时间分层，不作整体化推断。
\eventanchor{EQ-1685-ALBAZIN-FIRST-PRELUDE}{「清军围攻并迫使俄方撤离阿尔…}康熙二十四年（1685年），「清军围攻并迫使俄方撤离阿尔巴津，黑龙江边境进入直接谈判前的军事对峙」之前的条件、信息和争议已通过中央与地方材料显现，构成该事件可追踪的前史节点。核心取证：《清圣祖实录》康熙二十四年黑龙江军报；《清史稿·邦交志》；俄国西伯利亚衙门文件与尼布楚谈判材料；同时期地方档案、地方志、日记或对方材料应按具体地区补检。该项锚定的是前史条件与信息背景，不把条件直接写成唯一因果。来源保留：官修编年和地方材料的记录密度与立场并不相同；本行不将条件、传闻、呈报或后见解释直接等同为确定因果。
\eventanchor{EQ-1685-ALBAZIN-FIRST-DECISION}{围绕「清军围攻并迫使俄方撤离…}康熙二十四年（1685年），围绕「清军围攻并迫使俄方撤离阿尔巴津，黑龙江边境进入直接谈判前的军事对峙」的命令、调度、照会或呈报形成独立的中央—地方行动文书链。核心取证：《清圣祖实录》康熙二十四年黑龙江军报；《清史稿·邦交志》；俄国西伯利亚衙门文件与尼布楚谈判材料。该项锚定的是命令或决策环节，命令抵达和结果仍须另外核验。来源保留：奏折、上谕、部议、军机档或条约可证明行政动作和机构立场，不自动证明所有地区、群体或对手按其意图行动。
\eventanchor{EQ-1685-ALBAZIN-FIRST-AFTERMATH}{「清军围攻并迫使俄方撤离阿尔…}康熙二十四年（1685年），「清军围攻并迫使俄方撤离阿尔巴津，黑龙江边境进入直接谈判前的军事对峙」引发的善后、执行或地方回应构成可由不同材料追踪的后续节点。核心取证：《清圣祖实录》康熙二十四年黑龙江军报；《清史稿·邦交志》；俄国西伯利亚衙门文件与尼布楚谈判材料；相关地区的档案、志书、契约、报刊或对方材料。该项锚定的是后续追踪问题，影响范围须按地区与群体分别判断。来源保留：战后、改革后或条约后的记忆、地方记录和官方总结常具有事后选择性；后续影响须按地点、群体和时间分层，不作整体化推断。
\eventanchor{EQ-1686-ALBAZIN-SECOND-PRELUDE}{「俄方重返阿尔巴津后，清军再…}康熙二十五年（1686年），「俄方重返阿尔巴津后，清军再次围困，双方在疾病、补给与外交压力下僵持」之前的条件、信息和争议已通过中央与地方材料显现，构成该事件可追踪的前史节点。核心取证：《清圣祖实录》康熙二十五年黑龙江军报；《清史稿·邦交志》；俄国阿尔巴津驻军报告及城址考古报告；同时期地方档案、地方志、日记或对方材料应按具体地区补检。该项锚定的是前史条件与信息背景，不把条件直接写成唯一因果。来源保留：官修编年和地方材料的记录密度与立场并不相同；本行不将条件、传闻、呈报或后见解释直接等同为确定因果。
\eventanchor{EQ-1686-ALBAZIN-SECOND-DECISION}{围绕「俄方重返阿尔巴津后，清…}康熙二十五年（1686年），围绕「俄方重返阿尔巴津后，清军再次围困，双方在疾病、补给与外交压力下僵持」的命令、调度、照会或呈报形成独立的中央—地方行动文书链。核心取证：《清圣祖实录》康熙二十五年黑龙江军报；《清史稿·邦交志》；俄国阿尔巴津驻军报告及城址考古报告。该项锚定的是命令或决策环节，命令抵达和结果仍须另外核验。来源保留：奏折、上谕、部议、军机档或条约可证明行政动作和机构立场，不自动证明所有地区、群体或对手按其意图行动。
\eventanchor{EQ-1686-ALBAZIN-SECOND-AFTERMATH}{「俄方重返阿尔巴津后，清军再…}康熙二十五年（1686年），「俄方重返阿尔巴津后，清军再次围困，双方在疾病、补给与外交压力下僵持」引发的善后、执行或地方回应构成可由不同材料追踪的后续节点。核心取证：《清圣祖实录》康熙二十五年黑龙江军报；《清史稿·邦交志》；俄国阿尔巴津驻军报告及城址考古报告；相关地区的档案、志书、契约、报刊或对方材料。该项锚定的是后续追踪问题，影响范围须按地区与群体分别判断。来源保留：战后、改革后或条约后的记忆、地方记录和官方总结常具有事后选择性；后续影响须按地点、群体和时间分层，不作整体化推断。
\eventanchor{EQ-1688-GALDAN-KHALKHA-PRELUDE}{「噶尔丹进攻喀尔喀，喀尔喀诸…}康熙二十七年（1688年），「噶尔丹进攻喀尔喀，喀尔喀诸部南徙并向清廷求援，北方草原秩序改变」之前的条件、信息和争议已通过中央与地方材料显现，构成该事件可追踪的前史节点。核心取证：《清圣祖实录》康熙二十七年蒙古军报；《清史稿·噶尔丹传》；《蒙古源流》与俄罗斯使团材料；同时期地方档案、地方志、日记或对方材料应按具体地区补检。该项锚定的是前史条件与信息背景，不把条件直接写成唯一因果。来源保留：官修编年和地方材料的记录密度与立场并不相同；本行不将条件、传闻、呈报或后见解释直接等同为确定因果。
\eventanchor{EQ-1688-GALDAN-KHALKHA-DECISION}{围绕「噶尔丹进攻喀尔喀，喀尔…}康熙二十七年（1688年），围绕「噶尔丹进攻喀尔喀，喀尔喀诸部南徙并向清廷求援，北方草原秩序改变」的命令、调度、照会或呈报形成独立的中央—地方行动文书链。核心取证：《清圣祖实录》康熙二十七年蒙古军报；《清史稿·噶尔丹传》；《蒙古源流》与俄罗斯使团材料。该项锚定的是命令或决策环节，命令抵达和结果仍须另外核验。来源保留：奏折、上谕、部议、军机档或条约可证明行政动作和机构立场，不自动证明所有地区、群体或对手按其意图行动。
\eventanchor{EQ-1688-GALDAN-KHALKHA-AFTERMATH}{「噶尔丹进攻喀尔喀，喀尔喀诸…}康熙二十七年（1688年），「噶尔丹进攻喀尔喀，喀尔喀诸部南徙并向清廷求援，北方草原秩序改变」引发的善后、执行或地方回应构成可由不同材料追踪的后续节点。核心取证：《清圣祖实录》康熙二十七年蒙古军报；《清史稿·噶尔丹传》；《蒙古源流》与俄罗斯使团材料；相关地区的档案、志书、契约、报刊或对方材料。该项锚定的是后续追踪问题，影响范围须按地区与群体分别判断。来源保留：战后、改革后或条约后的记忆、地方记录和官方总结常具有事后选择性；后续影响须按地点、群体和时间分层，不作整体化推断。
\eventanchor{EQ-1689-NERCHINSK-PRELUDE}{「清俄在尼布楚签订条约，以满…}康熙二十八年七月（1689年），「清俄在尼布楚签订条约，以满文、拉丁文和俄文文本处理边界与阿尔巴津问题」之前的条件、信息和争议已通过中央与地方材料显现，构成该事件可追踪的前史节点。核心取证：《清圣祖实录》康熙二十八年七月；《清史稿·邦交志》；《尼布楚条约》满文、拉丁文、俄文文本及俄国谈判档；同时期地方档案、地方志、日记或对方材料应按具体地区补检。该项锚定的是前史条件与信息背景，不把条件直接写成唯一因果。来源保留：官修编年和地方材料的记录密度与立场并不相同；本行不将条件、传闻、呈报或后见解释直接等同为确定因果。
\eventanchor{EQ-1689-NERCHINSK-DECISION}{围绕「清俄在尼布楚签订条约，…}康熙二十八年七月（1689年），围绕「清俄在尼布楚签订条约，以满文、拉丁文和俄文文本处理边界与阿尔巴津问题」的命令、调度、照会或呈报形成独立的中央—地方行动文书链。核心取证：《清圣祖实录》康熙二十八年七月；《清史稿·邦交志》；《尼布楚条约》满文、拉丁文、俄文文本及俄国谈判档。该项锚定的是命令或决策环节，命令抵达和结果仍须另外核验。来源保留：奏折、上谕、部议、军机档或条约可证明行政动作和机构立场，不自动证明所有地区、群体或对手按其意图行动。
\eventanchor{EQ-1689-NERCHINSK-AFTERMATH}{「清俄在尼布楚签订条约，以满…}康熙二十八年七月（1689年），「清俄在尼布楚签订条约，以满文、拉丁文和俄文文本处理边界与阿尔巴津问题」引发的善后、执行或地方回应构成可由不同材料追踪的后续节点。核心取证：《清圣祖实录》康熙二十八年七月；《清史稿·邦交志》；《尼布楚条约》满文、拉丁文、俄文文本及俄国谈判档；相关地区的档案、志书、契约、报刊或对方材料。该项锚定的是后续追踪问题，影响范围须按地区与群体分别判断。来源保留：战后、改革后或条约后的记忆、地方记录和官方总结常具有事后选择性；后续影响须按地点、群体和时间分层，不作整体化推断。
\subsection{1690—1699}
\eventanchor{EQ-1690-ULAN-BUTUNG-PRELUDE}{「清军在乌兰布通与噶尔丹军交…}康熙二十九年（1690年），「清军在乌兰布通与噶尔丹军交战，阻止其向内蒙古腹地深入但未立即消灭其军队」之前的条件、信息和争议已通过中央与地方材料显现，构成该事件可追踪的前史节点。核心取证：《清圣祖实录》康熙二十九年八月；《清史稿·噶尔丹传》；《亲征平定朔漠方略》及俄国商旅报告；同时期地方档案、地方志、日记或对方材料应按具体地区补检。该项锚定的是前史条件与信息背景，不把条件直接写成唯一因果。来源保留：官修编年和地方材料的记录密度与立场并不相同；本行不将条件、传闻、呈报或后见解释直接等同为确定因果。
\eventanchor{EQ-1690-ULAN-BUTUNG-DECISION}{围绕「清军在乌兰布通与噶尔丹…}康熙二十九年（1690年），围绕「清军在乌兰布通与噶尔丹军交战，阻止其向内蒙古腹地深入但未立即消灭其军队」的命令、调度、照会或呈报形成独立的中央—地方行动文书链。核心取证：《清圣祖实录》康熙二十九年八月；《清史稿·噶尔丹传》；《亲征平定朔漠方略》及俄国商旅报告。该项锚定的是命令或决策环节，命令抵达和结果仍须另外核验。来源保留：奏折、上谕、部议、军机档或条约可证明行政动作和机构立场，不自动证明所有地区、群体或对手按其意图行动。
\eventanchor{EQ-1690-ULAN-BUTUNG-AFTERMATH}{「清军在乌兰布通与噶尔丹军交…}康熙二十九年（1690年），「清军在乌兰布通与噶尔丹军交战，阻止其向内蒙古腹地深入但未立即消灭其军队」引发的善后、执行或地方回应构成可由不同材料追踪的后续节点。核心取证：《清圣祖实录》康熙二十九年八月；《清史稿·噶尔丹传》；《亲征平定朔漠方略》及俄国商旅报告；相关地区的档案、志书、契约、报刊或对方材料。该项锚定的是后续追踪问题，影响范围须按地区与群体分别判断。来源保留：战后、改革后或条约后的记忆、地方记录和官方总结常具有事后选择性；后续影响须按地点、群体和时间分层，不作整体化推断。
\eventanchor{EQ-1691-DOLONNOR-PRELUDE}{「康熙帝在多伦诺尔会盟，安排…}康熙三十年（1691年），「康熙帝在多伦诺尔会盟，安排喀尔喀诸部入盟旗体系并调度对噶尔丹资源」之前的条件、信息和争议已通过中央与地方材料显现，构成该事件可追踪的前史节点。核心取证：《清圣祖实录》康熙三十年五月；《清史稿·藩部传》；理藩院满蒙文档案中的喀尔喀盟旗与赈济文书；同时期地方档案、地方志、日记或对方材料应按具体地区补检。该项锚定的是前史条件与信息背景，不把条件直接写成唯一因果。来源保留：官修编年和地方材料的记录密度与立场并不相同；本行不将条件、传闻、呈报或后见解释直接等同为确定因果。
\eventanchor{EQ-1691-DOLONNOR-DECISION}{围绕「康熙帝在多伦诺尔会盟，…}康熙三十年（1691年），围绕「康熙帝在多伦诺尔会盟，安排喀尔喀诸部入盟旗体系并调度对噶尔丹资源」的命令、调度、照会或呈报形成独立的中央—地方行动文书链。核心取证：《清圣祖实录》康熙三十年五月；《清史稿·藩部传》；理藩院满蒙文档案中的喀尔喀盟旗与赈济文书。该项锚定的是命令或决策环节，命令抵达和结果仍须另外核验。来源保留：奏折、上谕、部议、军机档或条约可证明行政动作和机构立场，不自动证明所有地区、群体或对手按其意图行动。
\eventanchor{EQ-1691-DOLONNOR-AFTERMATH}{「康熙帝在多伦诺尔会盟，安排…}康熙三十年（1691年），「康熙帝在多伦诺尔会盟，安排喀尔喀诸部入盟旗体系并调度对噶尔丹资源」引发的善后、执行或地方回应构成可由不同材料追踪的后续节点。核心取证：《清圣祖实录》康熙三十年五月；《清史稿·藩部传》；理藩院满蒙文档案中的喀尔喀盟旗与赈济文书；相关地区的档案、志书、契约、报刊或对方材料。该项锚定的是后续追踪问题，影响范围须按地区与群体分别判断。来源保留：战后、改革后或条约后的记忆、地方记录和官方总结常具有事后选择性；后续影响须按地点、群体和时间分层，不作整体化推断。
\eventanchor{EQ-1695-GALDAN-CAMPAIGN-PRELUDE}{「康熙帝筹划并发动对噶尔丹的…}康熙三十四年（1695年），「康熙帝筹划并发动对噶尔丹的第二阶段远征，集中军粮、驼马和多路部队」之前的条件、信息和争议已通过中央与地方材料显现，构成该事件可追踪的前史节点。核心取证：《清圣祖实录》康熙三十四年军务条；《清史稿·圣祖本纪》；《康熙朝满文朱批奏折全译》噶尔丹军需奏折；同时期地方档案、地方志、日记或对方材料应按具体地区补检。该项锚定的是前史条件与信息背景，不把条件直接写成唯一因果。来源保留：官修编年和地方材料的记录密度与立场并不相同；本行不将条件、传闻、呈报或后见解释直接等同为确定因果。
\eventanchor{EQ-1695-GALDAN-CAMPAIGN-DECISION}{围绕「康熙帝筹划并发动对噶尔…}康熙三十四年（1695年），围绕「康熙帝筹划并发动对噶尔丹的第二阶段远征，集中军粮、驼马和多路部队」的命令、调度、照会或呈报形成独立的中央—地方行动文书链。核心取证：《清圣祖实录》康熙三十四年军务条；《清史稿·圣祖本纪》；《康熙朝满文朱批奏折全译》噶尔丹军需奏折。该项锚定的是命令或决策环节，命令抵达和结果仍须另外核验。来源保留：奏折、上谕、部议、军机档或条约可证明行政动作和机构立场，不自动证明所有地区、群体或对手按其意图行动。
\eventanchor{EQ-1695-GALDAN-CAMPAIGN-AFTERMATH}{「康熙帝筹划并发动对噶尔丹的…}康熙三十四年（1695年），「康熙帝筹划并发动对噶尔丹的第二阶段远征，集中军粮、驼马和多路部队」引发的善后、执行或地方回应构成可由不同材料追踪的后续节点。核心取证：《清圣祖实录》康熙三十四年军务条；《清史稿·圣祖本纪》；《康熙朝满文朱批奏折全译》噶尔丹军需奏折；相关地区的档案、志书、契约、报刊或对方材料。该项锚定的是后续追踪问题，影响范围须按地区与群体分别判断。来源保留：战后、改革后或条约后的记忆、地方记录和官方总结常具有事后选择性；后续影响须按地点、群体和时间分层，不作整体化推断。
\eventanchor{EQ-1696-JAOMODO-PRELUDE}{「清军在昭莫多一带击败噶尔丹…}康熙三十五年（1696年），「清军在昭莫多一带击败噶尔丹主力，噶尔丹向西遁走」之前的条件、信息和争议已通过中央与地方材料显现，构成该事件可追踪的前史节点。核心取证：《清圣祖实录》康熙三十五年五月；《清史稿·噶尔丹传》；《亲征平定朔漠方略》与俄国边境报告；同时期地方档案、地方志、日记或对方材料应按具体地区补检。该项锚定的是前史条件与信息背景，不把条件直接写成唯一因果。来源保留：官修编年和地方材料的记录密度与立场并不相同；本行不将条件、传闻、呈报或后见解释直接等同为确定因果。
\eventanchor{EQ-1696-JAOMODO-DECISION}{围绕「清军在昭莫多一带击败噶…}康熙三十五年（1696年），围绕「清军在昭莫多一带击败噶尔丹主力，噶尔丹向西遁走」的命令、调度、照会或呈报形成独立的中央—地方行动文书链。核心取证：《清圣祖实录》康熙三十五年五月；《清史稿·噶尔丹传》；《亲征平定朔漠方略》与俄国边境报告。该项锚定的是命令或决策环节，命令抵达和结果仍须另外核验。来源保留：奏折、上谕、部议、军机档或条约可证明行政动作和机构立场，不自动证明所有地区、群体或对手按其意图行动。
\eventanchor{EQ-1696-JAOMODO-AFTERMATH}{「清军在昭莫多一带击败噶尔丹…}康熙三十五年（1696年），「清军在昭莫多一带击败噶尔丹主力，噶尔丹向西遁走」引发的善后、执行或地方回应构成可由不同材料追踪的后续节点。核心取证：《清圣祖实录》康熙三十五年五月；《清史稿·噶尔丹传》；《亲征平定朔漠方略》与俄国边境报告；相关地区的档案、志书、契约、报刊或对方材料。该项锚定的是后续追踪问题，影响范围须按地区与群体分别判断。来源保留：战后、改革后或条约后的记忆、地方记录和官方总结常具有事后选择性；后续影响须按地点、群体和时间分层，不作整体化推断。
\eventanchor{EQ-1697-GALDAN-DEATH-PRELUDE}{「噶尔丹在科布多一带去世，其…}康熙三十六年（1697年），「噶尔丹在科布多一带去世，其部众分散或向准噶尔核心地区退却」之前的条件、信息和争议已通过中央与地方材料显现，构成该事件可追踪的前史节点。核心取证：《清圣祖实录》康熙三十六年；《清史稿·噶尔丹传》；俄国西伯利亚档案与蒙古口传整理；同时期地方档案、地方志、日记或对方材料应按具体地区补检。该项锚定的是前史条件与信息背景，不把条件直接写成唯一因果。来源保留：官修编年和地方材料的记录密度与立场并不相同；本行不将条件、传闻、呈报或后见解释直接等同为确定因果。
\eventanchor{EQ-1697-GALDAN-DEATH-DECISION}{围绕「噶尔丹在科布多一带去世…}康熙三十六年（1697年），围绕「噶尔丹在科布多一带去世，其部众分散或向准噶尔核心地区退却」的命令、调度、照会或呈报形成独立的中央—地方行动文书链。核心取证：《清圣祖实录》康熙三十六年；《清史稿·噶尔丹传》；俄国西伯利亚档案与蒙古口传整理。该项锚定的是命令或决策环节，命令抵达和结果仍须另外核验。来源保留：奏折、上谕、部议、军机档或条约可证明行政动作和机构立场，不自动证明所有地区、群体或对手按其意图行动。
\eventanchor{EQ-1697-GALDAN-DEATH-AFTERMATH}{「噶尔丹在科布多一带去世，其…}康熙三十六年（1697年），「噶尔丹在科布多一带去世，其部众分散或向准噶尔核心地区退却」引发的善后、执行或地方回应构成可由不同材料追踪的后续节点。核心取证：《清圣祖实录》康熙三十六年；《清史稿·噶尔丹传》；俄国西伯利亚档案与蒙古口传整理；相关地区的档案、志书、契约、报刊或对方材料。该项锚定的是后续追踪问题，影响范围须按地区与群体分别判断。来源保留：战后、改革后或条约后的记忆、地方记录和官方总结常具有事后选择性；后续影响须按地点、群体和时间分层，不作整体化推断。
\eventanchor{EQ-1698-KHALKHA-ADMINISTRATION-PRELUDE}{「清廷在噶尔丹败亡后调整喀尔…}康熙三十七年（1698年），「清廷在噶尔丹败亡后调整喀尔喀盟旗、台站与赈济安排，将战时同盟转为日常治理」之前的条件、信息和争议已通过中央与地方材料显现，构成该事件可追踪的前史节点。核心取证：《清圣祖实录》康熙三十七年蒙古条；《清史稿·藩部传》；理藩院满蒙文档案中的喀尔喀盟旗与赈济文书；同时期地方档案、地方志、日记或对方材料应按具体地区补检。该项锚定的是前史条件与信息背景，不把条件直接写成唯一因果。来源保留：官修编年和地方材料的记录密度与立场并不相同；本行不将条件、传闻、呈报或后见解释直接等同为确定因果。
\eventanchor{EQ-1698-KHALKHA-ADMINISTRATION-DECISION}{围绕「清廷在噶尔丹败亡后调整…}康熙三十七年（1698年），围绕「清廷在噶尔丹败亡后调整喀尔喀盟旗、台站与赈济安排，将战时同盟转为日常治理」的命令、调度、照会或呈报形成独立的中央—地方行动文书链。核心取证：《清圣祖实录》康熙三十七年蒙古条；《清史稿·藩部传》；理藩院满蒙文档案中的喀尔喀盟旗与赈济文书。该项锚定的是命令或决策环节，命令抵达和结果仍须另外核验。来源保留：奏折、上谕、部议、军机档或条约可证明行政动作和机构立场，不自动证明所有地区、群体或对手按其意图行动。
\eventanchor{EQ-1698-KHALKHA-ADMINISTRATION-AFTERMATH}{「清廷在噶尔丹败亡后调整喀尔…}康熙三十七年（1698年），「清廷在噶尔丹败亡后调整喀尔喀盟旗、台站与赈济安排，将战时同盟转为日常治理」引发的善后、执行或地方回应构成可由不同材料追踪的后续节点。核心取证：《清圣祖实录》康熙三十七年蒙古条；《清史稿·藩部传》；理藩院满蒙文档案中的喀尔喀盟旗与赈济文书；相关地区的档案、志书、契约、报刊或对方材料。该项锚定的是后续追踪问题，影响范围须按地区与群体分别判断。来源保留：战后、改革后或条约后的记忆、地方记录和官方总结常具有事后选择性；后续影响须按地点、群体和时间分层，不作整体化推断。
\subsection{1700—1709}
\eventanchor{EQ-1705-LHAZANG-TAKEOVER-PRELUDE}{「拉藏汗推翻第巴桑结嘉措的政…}康熙四十四年（1705年），「拉藏汗推翻第巴桑结嘉措的政权安排，西藏政教权力结构剧烈变化」之前的条件、信息和争议已通过中央与地方材料显现，构成该事件可追踪的前史节点。核心取证：《清圣祖实录》康熙四十四年西藏条；《清史稿·西藏传》；《西藏王臣记》藏文与清廷驻藏来文；同时期地方档案、地方志、日记或对方材料应按具体地区补检。该项锚定的是前史条件与信息背景，不把条件直接写成唯一因果。来源保留：官修编年和地方材料的记录密度与立场并不相同；本行不将条件、传闻、呈报或后见解释直接等同为确定因果。
\eventanchor{EQ-1705-LHAZANG-TAKEOVER-DECISION}{围绕「拉藏汗推翻第巴桑结嘉措…}康熙四十四年（1705年），围绕「拉藏汗推翻第巴桑结嘉措的政权安排，西藏政教权力结构剧烈变化」的命令、调度、照会或呈报形成独立的中央—地方行动文书链。核心取证：《清圣祖实录》康熙四十四年西藏条；《清史稿·西藏传》；《西藏王臣记》藏文与清廷驻藏来文。该项锚定的是命令或决策环节，命令抵达和结果仍须另外核验。来源保留：奏折、上谕、部议、军机档或条约可证明行政动作和机构立场，不自动证明所有地区、群体或对手按其意图行动。
\eventanchor{EQ-1705-LHAZANG-TAKEOVER-AFTERMATH}{「拉藏汗推翻第巴桑结嘉措的政…}康熙四十四年（1705年），「拉藏汗推翻第巴桑结嘉措的政权安排，西藏政教权力结构剧烈变化」引发的善后、执行或地方回应构成可由不同材料追踪的后续节点。核心取证：《清圣祖实录》康熙四十四年西藏条；《清史稿·西藏传》；《西藏王臣记》藏文与清廷驻藏来文；相关地区的档案、志书、契约、报刊或对方材料。该项锚定的是后续追踪问题，影响范围须按地区与群体分别判断。来源保留：战后、改革后或条约后的记忆、地方记录和官方总结常具有事后选择性；后续影响须按地点、群体和时间分层，不作整体化推断。
\eventanchor{EQ-1706-SIXTH-DALAI-PRELUDE}{「第六世达赖喇嘛仓央嘉措被押…}康熙四十五年（1706年），「第六世达赖喇嘛仓央嘉措被押离拉萨，途中死亡或失踪，其结局与继承合法性成为争议」之前的条件、信息和争议已通过中央与地方材料显现，构成该事件可追踪的前史节点。核心取证：《清圣祖实录》康熙四十五年西藏条；《清史稿·西藏传》；第巴传记、藏文史书与拉达克、蒙古相关记载；同时期地方档案、地方志、日记或对方材料应按具体地区补检。该项锚定的是前史条件与信息背景，不把条件直接写成唯一因果。来源保留：官修编年和地方材料的记录密度与立场并不相同；本行不将条件、传闻、呈报或后见解释直接等同为确定因果。
\eventanchor{EQ-1706-SIXTH-DALAI-DECISION}{围绕「第六世达赖喇嘛仓央嘉措…}康熙四十五年（1706年），围绕「第六世达赖喇嘛仓央嘉措被押离拉萨，途中死亡或失踪，其结局与继承合法性成为争议」的命令、调度、照会或呈报形成独立的中央—地方行动文书链。核心取证：《清圣祖实录》康熙四十五年西藏条；《清史稿·西藏传》；第巴传记、藏文史书与拉达克、蒙古相关记载。该项锚定的是命令或决策环节，命令抵达和结果仍须另外核验。来源保留：奏折、上谕、部议、军机档或条约可证明行政动作和机构立场，不自动证明所有地区、群体或对手按其意图行动。
\eventanchor{EQ-1706-SIXTH-DALAI-AFTERMATH}{「第六世达赖喇嘛仓央嘉措被押…}康熙四十五年（1706年），「第六世达赖喇嘛仓央嘉措被押离拉萨，途中死亡或失踪，其结局与继承合法性成为争议」引发的善后、执行或地方回应构成可由不同材料追踪的后续节点。核心取证：《清圣祖实录》康熙四十五年西藏条；《清史稿·西藏传》；第巴传记、藏文史书与拉达克、蒙古相关记载；相关地区的档案、志书、契约、报刊或对方材料。该项锚定的是后续追踪问题，影响范围须按地区与群体分别判断。来源保留：战后、改革后或条约后的记忆、地方记录和官方总结常具有事后选择性；后续影响须按地点、群体和时间分层，不作整体化推断。
\eventanchor{EQ-1708-KELZANG-PRELUDE}{「卡尔藏嘉措在青海塔尔寺附近…}康熙四十七年（1708年），「卡尔藏嘉措在青海塔尔寺附近被认定为第六世达赖喇嘛转世之一，围绕其身份形成跨区域支持网络」之前的条件、信息和争议已通过中央与地方材料显现，构成该事件可追踪的前史节点。核心取证：《清圣祖实录》康熙四十七年青海西藏条；《清史稿·西藏传》；塔尔寺藏文档案与第七世达赖喇嘛传；同时期地方档案、地方志、日记或对方材料应按具体地区补检。该项锚定的是前史条件与信息背景，不把条件直接写成唯一因果。来源保留：官修编年和地方材料的记录密度与立场并不相同；本行不将条件、传闻、呈报或后见解释直接等同为确定因果。
\eventanchor{EQ-1708-KELZANG-DECISION}{围绕「卡尔藏嘉措在青海塔尔寺…}康熙四十七年（1708年），围绕「卡尔藏嘉措在青海塔尔寺附近被认定为第六世达赖喇嘛转世之一，围绕其身份形成跨区域支持网络」的命令、调度、照会或呈报形成独立的中央—地方行动文书链。核心取证：《清圣祖实录》康熙四十七年青海西藏条；《清史稿·西藏传》；塔尔寺藏文档案与第七世达赖喇嘛传。该项锚定的是命令或决策环节，命令抵达和结果仍须另外核验。来源保留：奏折、上谕、部议、军机档或条约可证明行政动作和机构立场，不自动证明所有地区、群体或对手按其意图行动。
\eventanchor{EQ-1708-KELZANG-AFTERMATH}{「卡尔藏嘉措在青海塔尔寺附近…}康熙四十七年（1708年），「卡尔藏嘉措在青海塔尔寺附近被认定为第六世达赖喇嘛转世之一，围绕其身份形成跨区域支持网络」引发的善后、执行或地方回应构成可由不同材料追踪的后续节点。核心取证：《清圣祖实录》康熙四十七年青海西藏条；《清史稿·西藏传》；塔尔寺藏文档案与第七世达赖喇嘛传；相关地区的档案、志书、契约、报刊或对方材料。该项锚定的是后续追踪问题，影响范围须按地区与群体分别判断。来源保留：战后、改革后或条约后的记忆、地方记录和官方总结常具有事后选择性；后续影响须按地点、群体和时间分层，不作整体化推断。
\subsection{1710—1719}
\eventanchor{EQ-1711-TAX-FREEZE-PRELUDE}{「康熙帝宣布以康熙五十年丁数…}康熙五十年（1711年），「康熙帝宣布以康熙五十年丁数为定额，后生人丁不再加征丁税」之前的条件、信息和争议已通过中央与地方材料显现，构成该事件可追踪的前史节点。核心取证：《清圣祖实录》康熙五十年丁银谕；《清史稿·食货志》；户部档案与《大清会典事例》赋役条；同时期地方档案、地方志、日记或对方材料应按具体地区补检。该项锚定的是前史条件与信息背景，不把条件直接写成唯一因果。来源保留：官修编年和地方材料的记录密度与立场并不相同；本行不将条件、传闻、呈报或后见解释直接等同为确定因果。
\eventanchor{EQ-1711-TAX-FREEZE-DECISION}{围绕「康熙帝宣布以康熙五十年…}康熙五十年（1711年），围绕「康熙帝宣布以康熙五十年丁数为定额，后生人丁不再加征丁税」的命令、调度、照会或呈报形成独立的中央—地方行动文书链。核心取证：《清圣祖实录》康熙五十年丁银谕；《清史稿·食货志》；户部档案与《大清会典事例》赋役条。该项锚定的是命令或决策环节，命令抵达和结果仍须另外核验。来源保留：奏折、上谕、部议、军机档或条约可证明行政动作和机构立场，不自动证明所有地区、群体或对手按其意图行动。
\eventanchor{EQ-1711-TAX-FREEZE-AFTERMATH}{「康熙帝宣布以康熙五十年丁数…}康熙五十年（1711年），「康熙帝宣布以康熙五十年丁数为定额，后生人丁不再加征丁税」引发的善后、执行或地方回应构成可由不同材料追踪的后续节点。核心取证：《清圣祖实录》康熙五十年丁银谕；《清史稿·食货志》；户部档案与《大清会典事例》赋役条；相关地区的档案、志书、契约、报刊或对方材料。该项锚定的是后续追踪问题，影响范围须按地区与群体分别判断。来源保留：战后、改革后或条约后的记忆、地方记录和官方总结常具有事后选择性；后续影响须按地点、群体和时间分层，不作整体化推断。
\eventanchor{EQ-1712-TAX-EDICT-PRELUDE}{「清廷再次谕令固定丁额并要求…}康熙五十一年（1712年），「清廷再次谕令固定丁额并要求地方遵行，围绕丁银征收与造册形成持续行政问题」之前的条件、信息和争议已通过中央与地方材料显现，构成该事件可追踪的前史节点。核心取证：《清圣祖实录》康熙五十一年赋役谕；《清史稿·食货志》；户部题本与江苏、直隶等省赋役档；同时期地方档案、地方志、日记或对方材料应按具体地区补检。该项锚定的是前史条件与信息背景，不把条件直接写成唯一因果。来源保留：官修编年和地方材料的记录密度与立场并不相同；本行不将条件、传闻、呈报或后见解释直接等同为确定因果。
\eventanchor{EQ-1712-TAX-EDICT-DECISION}{围绕「清廷再次谕令固定丁额并…}康熙五十一年（1712年），围绕「清廷再次谕令固定丁额并要求地方遵行，围绕丁银征收与造册形成持续行政问题」的命令、调度、照会或呈报形成独立的中央—地方行动文书链。核心取证：《清圣祖实录》康熙五十一年赋役谕；《清史稿·食货志》；户部题本与江苏、直隶等省赋役档。该项锚定的是命令或决策环节，命令抵达和结果仍须另外核验。来源保留：奏折、上谕、部议、军机档或条约可证明行政动作和机构立场，不自动证明所有地区、群体或对手按其意图行动。
\eventanchor{EQ-1712-TAX-EDICT-AFTERMATH}{「清廷再次谕令固定丁额并要求…}康熙五十一年（1712年），「清廷再次谕令固定丁额并要求地方遵行，围绕丁银征收与造册形成持续行政问题」引发的善后、执行或地方回应构成可由不同材料追踪的后续节点。核心取证：《清圣祖实录》康熙五十一年赋役谕；《清史稿·食货志》；户部题本与江苏、直隶等省赋役档；相关地区的档案、志书、契约、报刊或对方材料。该项锚定的是后续追踪问题，影响范围须按地区与群体分别判断。来源保留：战后、改革后或条约后的记忆、地方记录和官方总结常具有事后选择性；后续影响须按地点、群体和时间分层，不作整体化推断。
\eventanchor{EQ-1717-DZUNGAR-LHASA-PRELUDE}{「准噶尔军进入拉萨并杀拉藏汗…}康熙五十六年（1717年），「准噶尔军进入拉萨并杀拉藏汗，拉萨政局再次易手」之前的条件、信息和争议已通过中央与地方材料显现，构成该事件可追踪的前史节点。核心取证：《清圣祖实录》康熙五十六年西藏军报；《清史稿·西藏传》；藏文《西藏王臣记》续编与准噶尔相关材料；同时期地方档案、地方志、日记或对方材料应按具体地区补检。该项锚定的是前史条件与信息背景，不把条件直接写成唯一因果。来源保留：官修编年和地方材料的记录密度与立场并不相同；本行不将条件、传闻、呈报或后见解释直接等同为确定因果。
\eventanchor{EQ-1717-DZUNGAR-LHASA-DECISION}{围绕「准噶尔军进入拉萨并杀拉…}康熙五十六年（1717年），围绕「准噶尔军进入拉萨并杀拉藏汗，拉萨政局再次易手」的命令、调度、照会或呈报形成独立的中央—地方行动文书链。核心取证：《清圣祖实录》康熙五十六年西藏军报；《清史稿·西藏传》；藏文《西藏王臣记》续编与准噶尔相关材料。该项锚定的是命令或决策环节，命令抵达和结果仍须另外核验。来源保留：奏折、上谕、部议、军机档或条约可证明行政动作和机构立场，不自动证明所有地区、群体或对手按其意图行动。
\eventanchor{EQ-1717-DZUNGAR-LHASA-AFTERMATH}{「准噶尔军进入拉萨并杀拉藏汗…}康熙五十六年（1717年），「准噶尔军进入拉萨并杀拉藏汗，拉萨政局再次易手」引发的善后、执行或地方回应构成可由不同材料追踪的后续节点。核心取证：《清圣祖实录》康熙五十六年西藏军报；《清史稿·西藏传》；藏文《西藏王臣记》续编与准噶尔相关材料；相关地区的档案、志书、契约、报刊或对方材料。该项锚定的是后续追踪问题，影响范围须按地区与群体分别判断。来源保留：战后、改革后或条约后的记忆、地方记录和官方总结常具有事后选择性；后续影响须按地点、群体和时间分层，不作整体化推断。
\eventanchor{EQ-1718-QING-TIBET-EXPEDITION-PRELUDE}{「清军自四川方向入藏，在喀喇…}康熙五十七年（1718年），「清军自四川方向入藏，在喀喇乌苏一带覆败，首批救援行动受挫」之前的条件、信息和争议已通过中央与地方材料显现，构成该事件可追踪的前史节点。核心取证：《清圣祖实录》康熙五十七年川藏军报；《清史稿·西藏传》；《康熙朝满文朱批奏折全译》入藏军需奏折及藏文地方记载；同时期地方档案、地方志、日记或对方材料应按具体地区补检。该项锚定的是前史条件与信息背景，不把条件直接写成唯一因果。来源保留：官修编年和地方材料的记录密度与立场并不相同；本行不将条件、传闻、呈报或后见解释直接等同为确定因果。
\eventanchor{EQ-1718-QING-TIBET-EXPEDITION-DECISION}{围绕「清军自四川方向入藏，在…}康熙五十七年（1718年），围绕「清军自四川方向入藏，在喀喇乌苏一带覆败，首批救援行动受挫」的命令、调度、照会或呈报形成独立的中央—地方行动文书链。核心取证：《清圣祖实录》康熙五十七年川藏军报；《清史稿·西藏传》；《康熙朝满文朱批奏折全译》入藏军需奏折及藏文地方记载。该项锚定的是命令或决策环节，命令抵达和结果仍须另外核验。来源保留：奏折、上谕、部议、军机档或条约可证明行政动作和机构立场，不自动证明所有地区、群体或对手按其意图行动。
\eventanchor{EQ-1718-QING-TIBET-EXPEDITION-AFTERMATH}{「清军自四川方向入藏，在喀喇…}康熙五十七年（1718年），「清军自四川方向入藏，在喀喇乌苏一带覆败，首批救援行动受挫」引发的善后、执行或地方回应构成可由不同材料追踪的后续节点。核心取证：《清圣祖实录》康熙五十七年川藏军报；《清史稿·西藏传》；《康熙朝满文朱批奏折全译》入藏军需奏折及藏文地方记载；相关地区的档案、志书、契约、报刊或对方材料。该项锚定的是后续追踪问题，影响范围须按地区与群体分别判断。来源保留：战后、改革后或条约后的记忆、地方记录和官方总结常具有事后选择性；后续影响须按地点、群体和时间分层，不作整体化推断。
\subsection{1720—1729}
\eventanchor{EQ-1720-LHASA-RESTORATION-PRELUDE}{「清军及青海蒙古盟军进入拉萨…}康熙五十九年（1720年），「清军及青海蒙古盟军进入拉萨，驱逐准噶尔力量并迎立卡尔藏嘉措为第七世达赖喇嘛」之前的条件、信息和争议已通过中央与地方材料显现，构成该事件可追踪的前史节点。核心取证：《清圣祖实录》康熙五十九年西藏军报；《清史稿·西藏传》；满文军机前身档、藏文第七世达赖传与青海蒙古文书；同时期地方档案、地方志、日记或对方材料应按具体地区补检。该项锚定的是前史条件与信息背景，不把条件直接写成唯一因果。来源保留：官修编年和地方材料的记录密度与立场并不相同；本行不将条件、传闻、呈报或后见解释直接等同为确定因果。
\eventanchor{EQ-1720-LHASA-RESTORATION-DECISION}{围绕「清军及青海蒙古盟军进入…}康熙五十九年（1720年），围绕「清军及青海蒙古盟军进入拉萨，驱逐准噶尔力量并迎立卡尔藏嘉措为第七世达赖喇嘛」的命令、调度、照会或呈报形成独立的中央—地方行动文书链。核心取证：《清圣祖实录》康熙五十九年西藏军报；《清史稿·西藏传》；满文军机前身档、藏文第七世达赖传与青海蒙古文书。该项锚定的是命令或决策环节，命令抵达和结果仍须另外核验。来源保留：奏折、上谕、部议、军机档或条约可证明行政动作和机构立场，不自动证明所有地区、群体或对手按其意图行动。
\eventanchor{EQ-1720-LHASA-RESTORATION-AFTERMATH}{「清军及青海蒙古盟军进入拉萨…}康熙五十九年（1720年），「清军及青海蒙古盟军进入拉萨，驱逐准噶尔力量并迎立卡尔藏嘉措为第七世达赖喇嘛」引发的善后、执行或地方回应构成可由不同材料追踪的后续节点。核心取证：《清圣祖实录》康熙五十九年西藏军报；《清史稿·西藏传》；满文军机前身档、藏文第七世达赖传与青海蒙古文书；相关地区的档案、志书、契约、报刊或对方材料。该项锚定的是后续追踪问题，影响范围须按地区与群体分别判断。来源保留：战后、改革后或条约后的记忆、地方记录和官方总结常具有事后选择性；后续影响须按地点、群体和时间分层，不作整体化推断。
\eventanchor{EQ-1721-ZHU-YIGUI-PRELUDE}{「朱一贵在台湾发动反清起事，…}康熙六十年（1721年），「朱一贵在台湾发动反清起事，一度攻入府城，清廷随后派兵恢复控制」之前的条件、信息和争议已通过中央与地方材料显现，构成该事件可追踪的前史节点。核心取证：《清圣祖实录》康熙六十年台湾军报；《清史稿·朱一贵传》；《台湾府志》、福建督抚奏折与契约文书；同时期地方档案、地方志、日记或对方材料应按具体地区补检。该项锚定的是前史条件与信息背景，不把条件直接写成唯一因果。来源保留：官修编年和地方材料的记录密度与立场并不相同；本行不将条件、传闻、呈报或后见解释直接等同为确定因果。
\eventanchor{EQ-1721-ZHU-YIGUI-DECISION}{围绕「朱一贵在台湾发动反清起…}康熙六十年（1721年），围绕「朱一贵在台湾发动反清起事，一度攻入府城，清廷随后派兵恢复控制」的命令、调度、照会或呈报形成独立的中央—地方行动文书链。核心取证：《清圣祖实录》康熙六十年台湾军报；《清史稿·朱一贵传》；《台湾府志》、福建督抚奏折与契约文书。该项锚定的是命令或决策环节，命令抵达和结果仍须另外核验。来源保留：奏折、上谕、部议、军机档或条约可证明行政动作和机构立场，不自动证明所有地区、群体或对手按其意图行动。
\eventanchor{EQ-1721-ZHU-YIGUI-AFTERMATH}{「朱一贵在台湾发动反清起事，…}康熙六十年（1721年），「朱一贵在台湾发动反清起事，一度攻入府城，清廷随后派兵恢复控制」引发的善后、执行或地方回应构成可由不同材料追踪的后续节点。核心取证：《清圣祖实录》康熙六十年台湾军报；《清史稿·朱一贵传》；《台湾府志》、福建督抚奏折与契约文书；相关地区的档案、志书、契约、报刊或对方材料。该项锚定的是后续追踪问题，影响范围须按地区与群体分别判断。来源保留：战后、改革后或条约后的记忆、地方记录和官方总结常具有事后选择性；后续影响须按地点、群体和时间分层，不作整体化推断。
\eventanchor{EQ-1722-KANGXI-DEATH-PRELUDE}{「康熙帝去世，皇四子胤禛即位…}康熙六十一年十一月（1722年），「康熙帝去世，皇四子胤禛即位为雍正帝，早期清朝的长期统治转入新的继承阶段」之前的条件、信息和争议已通过中央与地方材料显现，构成该事件可追踪的前史节点。核心取证：《清圣祖实录》康熙六十一年十一月；《清史稿·圣祖本纪、世宗本纪》；《雍正朝起居注册》与内阁大库遗诏档；同时期地方档案、地方志、日记或对方材料应按具体地区补检。该项锚定的是前史条件与信息背景，不把条件直接写成唯一因果。来源保留：官修编年和地方材料的记录密度与立场并不相同；本行不将条件、传闻、呈报或后见解释直接等同为确定因果。
\eventanchor{EQ-1722-KANGXI-DEATH-DECISION}{围绕「康熙帝去世，皇四子胤禛…}康熙六十一年十一月（1722年），围绕「康熙帝去世，皇四子胤禛即位为雍正帝，早期清朝的长期统治转入新的继承阶段」的命令、调度、照会或呈报形成独立的中央—地方行动文书链。核心取证：《清圣祖实录》康熙六十一年十一月；《清史稿·圣祖本纪、世宗本纪》；《雍正朝起居注册》与内阁大库遗诏档。该项锚定的是命令或决策环节，命令抵达和结果仍须另外核验。来源保留：奏折、上谕、部议、军机档或条约可证明行政动作和机构立场，不自动证明所有地区、群体或对手按其意图行动。
\eventanchor{EQ-1722-KANGXI-DEATH-AFTERMATH}{「康熙帝去世，皇四子胤禛即位…}康熙六十一年十一月（1722年），「康熙帝去世，皇四子胤禛即位为雍正帝，早期清朝的长期统治转入新的继承阶段」引发的善后、执行或地方回应构成可由不同材料追踪的后续节点。核心取证：《清圣祖实录》康熙六十一年十一月；《清史稿·圣祖本纪、世宗本纪》；《雍正朝起居注册》与内阁大库遗诏档；相关地区的档案、志书、契约、报刊或对方材料。该项锚定的是后续追踪问题，影响范围须按地区与群体分别判断。来源保留：战后、改革后或条约后的记忆、地方记录和官方总结常具有事后选择性；后续影响须按地点、群体和时间分层，不作整体化推断。
